\documentclass[12pt,a4paper,UTF8]{ctexart}

%==========================================================================
% 必要宏包
%==========================================================================
\usepackage{geometry}
\usepackage{amsmath}
\usepackage{amssymb}
\usepackage{graphicx}
\usepackage{booktabs}
\usepackage{hyperref}
\usepackage{float}
\usepackage{enumitem}
\usepackage{multirow}
\usepackage{array}
\usepackage{cite}
\usepackage{caption}

%==========================================================================
% 页面布局
%==========================================================================
\geometry{left=3cm,right=2.5cm,top=2.5cm,bottom=2.5cm}

%==========================================================================
% 超链接样式
%==========================================================================
\hypersetup{
    colorlinks=true,
    linkcolor=blue,
    citecolor=blue,
    urlcolor=blue
}

%==========================================================================
% 论文信息
%==========================================================================
\title{\textbf{基于耦合多智能体强化学习的\\低空航空交通协同管控模型}}
\author{彭义森\quad 指导老师：韩云祥\\[4pt]
\small 四川大学}
\date{2026年7月}

\begin{document}

\maketitle

%==========================================================================
% 摘要
%==========================================================================
\begin{abstract}
随着先进空中交通（AAM）的快速发展，电动垂直起降飞行器（eVTOL）在城市低空的高密度运行对传统空中交通管理方法提出了严峻挑战。现有多智能体强化学习（MARL）方法在时空耦合建模、稀疏奖励处理与环境不确定性建模三个维度仍存在不足。针对城市低空十字交叉路口的飞行器冲突避免与流量协同优化问题，本文提出一种耦合MARL低空交通管控模型。模型核心设计包括：（1）动态图注意力网络（DyGAT），融合空间注意力、时序注意力与基于物理量的连续动态边权重，刻画飞行器间时空耦合关联；（2）分层奖励函数，兼顾飞行安全与运行效率；（3）课程学习与基于冲突事件的轨迹加权训练策略（$\lambda=2.0$），在保持PPO on-policy理论框架的前提下提升稀疏冲突样本的利用效率。在融合风扰与传感器噪声的21架eVTOL交叉路口仿真场景中，本文方法相较固定轮询调度（RR）规则基线，冲突次数降低57.2\%，平均运行延迟缩短63.9\%；相较当前最优的CTDE-MARL方法MAPPO，冲突次数降低10.3\%（$p=0.180$，未达统计显著），但无冲突完成率（CFCR）从57.8\%提升至68.4\%（$p<0.001$，Cohen's $d \approx 5.27$），平均运行延迟缩短12.3\%。本文坦诚讨论了当前二维简化场景、均匀风扰模型、有限基线覆盖范围及小样本量（$n=5$）统计功效不足等局限性，并明确了投稿前的优先补充实验与长期研究方向。

\noindent\textbf{关键词：}低空航空交通管理；多智能体强化学习；动态图注意力网络；冲突避免；集中训练分布式执行；时空耦合建模；课程学习
\end{abstract}

%==========================================================================
% 1. 引言
%==========================================================================
\section{引言}

城市地面交通拥堵问题的日益加剧，推动了以eVTOL为核心的城市空中交通（UAM）产业快速发展。美国联邦航空管理局（FAA）预测，到2040年将有数百万架无人机与eVTOL在美国低空空域运行\cite{ref26}。与传统民航航路的结构化交通不同，城市低空交通具有高密度、强动态、高异构性三大特征——飞行器在三维空间中快速移动，邻域拓扑以秒级频率变化，不同制造商与载荷等级的飞行器具有差异化的飞行包线。在此背景下，交叉路口与汇合点的协同管控与冲突避免已成为低空无人机交通管理系统（UTM）面临的核心挑战。

传统空中交通流量管理（ATFM）多采用集中式优化架构，包括整数规划、蒙特卡洛仿真与基于交通流欧拉模型的线性规划等方法\cite{ref31}。这些方法依赖预先规划的飞行计划与静态航路设计，在可预测场景中效果良好，但计算复杂度随飞行器数量指数增长，难以适配低空环境中的实时动态不确定性。模型预测控制（MPC）\cite{ref27}、速度障碍法（ORCA/RVO）\cite{ref22}与控制屏障函数（CBF）\cite{ref23}等工程方法在形式化安全保证方面具有独特优势，但其性能高度依赖精确的环境动力学模型，且在高密度交互场景中倾向于保守。近年来，MARL因其分布式决策与自适应学习特性，被成功引入交通管控领域。

Agogino与Tumer的系列工作率先为ATFM建立了多智能体学习范式：从分布式航路点智能体调控框架\cite{ref1}，到基于差异奖励的大规模流量管理\cite{ref2}，系统性地验证了MARL在航空流量管理中的有效性。Brittain与Wei设计了基于A2C网络结合PPO损失函数的深度MARL框架，在BlueSky仿真器中实现了99.97\%的冲突规避率\cite{ref3}。Ghosh等提出深度集成MARL方法，通过元策略仲裁不同专家策略以提升决策鲁棒性\cite{ref9}。Xie等将强化学习与遗传算法结合，构建了面向UAM需求-容量平衡的混合学习框架\cite{ref6}。Li等最近提出的MS-GRL采用多尺度图增强强化学习解决密集无人机网络中的冲突解决问题，在100架UAV规模下验证了图结构对冲突规避的关键作用\cite{ref24}。

尽管上述工作取得了显著进展，现有研究在三个核心维度上仍存在不足：

\textbf{（1）时空耦合建模不足。}多数方法通过全连接网络或静态图结构编码智能体间关联。Brittain与Wei仅考虑最近邻静态关联\cite{ref3}；Spatharis等通过预定义关系矩阵编码交互\cite{ref7}；Ghosh等依赖网格化Geo-hash密度图\cite{ref9}；MS-GRL虽采用了多尺度图注意力机制，但其图结构基于固定距离阈值而非物理量驱动的连续动态权重\cite{ref24}。这些方法的共同局限在于：当飞行器快速移动导致邻域拓扑剧烈变化时，静态或半静态图结构无法及时反映飞行器间动态变化的冲突风险等级——两架相向飞行的飞行器在距离相同时的风险远高于同向飞行，但固定阈值图无法区分这两种情形。

\textbf{（2）稀疏奖励处理不足。}低空交通中冲突事件属于低概率高风险事件——在典型训练中，冲突样本占全部时间步的比例通常低于0.1\%。Aslani等指出标准$\epsilon$-greedy探索策略在高维连续动作空间中难以高效覆盖稀疏关键状态区域\cite{ref10}；翟子洋等的综述将稀疏奖励列为MARL交通信号控制的首要挑战\cite{ref11}。标准经验回放机制中冲突样本被淹没在海量正常样本中，且冲突规避的奖励信号仅在冲突发生或成功通过时才出现，导致信用分配困难——飞行器难以判断哪一时刻的速度调整最终规避了冲突。

\textbf{（3）环境不确定性建模不足。}多数研究采用简化模型，忽略低空环境的复杂性与不确定性。Spatharis等指出传统固定安全距离无法适配不同飞行速度下的制动需求\cite{ref7}。大多数现有研究将飞行器视为质点，忽略最大加速度、最大转弯角等物理运动约束。此外，低空风场建模多采用均匀随机扰动，与真实Dryden风湍流谱（MIL-F-8785C）\cite{ref20}或Kaimal谱模型（IEC 61400-1）的城市峡谷非均匀风场存在显著简化。需要诚实指出，本文当前的风扰模型同样采用均匀随机分布——这一简化对实验结果的影响在第5.4节中进行了坦诚讨论。

针对上述三个维度的不足，本文设计了一种耦合多智能体强化学习管控模型，核心思路为：（1）通过DyGAT的物理量驱动动态边权重精准刻画飞行器间时空耦合关系；（2）通过课程学习与基于冲突事件的episode级轨迹加权（$\lambda=2.0$），在严格保持PPO on-policy理论框架的前提下高效利用稀疏冲突样本；（3）通过动态安全距离与物理约束建模提升环境的工程合理性。

本文的主要贡献包括：

\begin{enumerate}[label=(\arabic*)]
    \item 构建了面向低空交叉路口冲突规避的仿真环境\texttt{AirTrafficEnv}，引入与速度正相关的动态安全距离模型与eVTOL物理运动约束（参数取值为参考当前主流eVTOL机型公开技术指标的代表性设定，详见第3.2.2节），并系统分析了风扰强度与传感器噪声对系统性能的影响；
    \item 提出基于DyGAT的耦合智能体群网络，通过距离、速度、航向角三个物理量直接计算连续动态边权重——消融实验表明该模块的移除使冲突次数增加173.6\%，是性能贡献最大的单一组件；
    \item 设计课程学习与基于冲突事件的episode级轨迹加权训练框架（$\lambda=2.0$），在不引入off-policy数据、不破坏PPO理论假设的前提下，使冲突样本的损失贡献翻倍——消融实验表明该模块的移除使冲突次数增加35.2\%；
    \item 构建双层评价指标体系——区分"表层指标"（累计奖励、总冲突次数、平均延迟）与"深层指标"（无冲突完成率CFCR、每架次冲突次数MCAE、最小间隔违反率MSVR），揭示PCR（99.1\%）与CFCR（68.4\%）之间的巨大差距，并引入Welch $t$检验诚实呈现TC不显著（$p=0.180$）与CFCR高度显著（$p<0.001$）的差异化统计图景；
    \item 与8种基线方法（涵盖传统规则、单智能体RL、独立多智能体学习与CTDE-MARL四类）进行系统对比，并坦诚讨论了MPC、ORCA/RVO、CBF等工程安全关键方法未纳入实验比较的局限。
\end{enumerate}

%==========================================================================
% 2. 相关工作
%==========================================================================
\section{相关工作}

本节从传统ATFM方法、MARL在交通管控中的应用及现有工作的差异化定位三个维度梳理相关研究，表\ref{tab:related_work}给出系统性对比。

\subsection{传统空中交通流量管理方法}

传统ATFM方法可分为集中式优化与分布式启发式两类。

\textbf{集中式优化方面，}Bertsimas等采用整数规划求解空域流量分配问题；Menon等通过交通流欧拉模型的线性规划实现动态调控\cite{ref31}。此类方法在可预测场景中效果良好，能够提供最优性保证，但在高密度、强动态场景中计算复杂度随问题规模指数增长，难以满足实时决策需求。近年来，基于混合整数线性规划（MILP）的集中式时空资源分配方法在中小规模场景中展现了良好的性能，但同样面临可扩展性挑战。

\textbf{启发式方法方面，}Agogino与Tumer设计了基于规则的反馈控制机制\cite{ref4}，计算效率高（每步决策时间$<1$ms），但性能上限受限于人工规则的表达能力，泛化至未见场景的能力较弱。MPC\cite{ref27}通过滚动时域优化实现对未来轨迹的预测性控制，在航空器间隔保证领域已有成熟应用；速度障碍法（ORCA/RVO）\cite{ref22}通过速度空间的几何约束实现分布式碰撞避免，计算效率高且具有形式化安全保证；CBF\cite{ref23}作为安全关键系统的形式化方法，通过李雅普诺夫类约束将安全需求编码为控制输入的限制条件，近年来在机器人导航与自动驾驶领域受到广泛关注。需要坦诚指出：受限于当前仿真环境的二维简化设定与资源约束，本文尚未将MPC、ORCA/RVO和CBF纳入实验对比——这是当前工作的一项重要局限，详见第5.4节的讨论。

\subsection{强化学习在交通管控中的应用}

\textbf{单智能体RL方面，}DQN已被应用于单飞行器冲突规避与单交叉口信号控制等场景。Aslani等在城市地面交通信号控制中系统对比了DQN与A2C在不同交通扰动事件下的鲁棒性\cite{ref10}。国内研究者在该方向也开展了大量工作：王迎\cite{ref13}、赖建辉\cite{ref16}、张新武\cite{ref17}分别在硕士论文中将深度强化学习应用于交通信号控制优化；赵晓华等\cite{ref14}、承向军等\cite{ref15}、卢守峰等\cite{ref18}探索了Q-learning及其变体在交叉口信号控制中的应用；江波等将DQN应用于机场道口协同智能调速\cite{ref19}；张春阳与席雅琪综述了人工智能在城市交通控制中的应用现状\cite{ref8}。需要说明，地面交通信号控制中的MARL方法在稀疏奖励处理、信用分配与CTDE架构等方面的设计思路对低空交通管控具有方法论层面的参考价值——两者均面临大规模协同决策与安全约束的核心挑战。

\textbf{多智能体RL领域，}Ghavamzadeh等提出层级式MARL框架，通过任务分解降低学习复杂度\cite{ref12}；Tumer与Agogino建立了基于航路点智能体的分布式ATM学习框架\cite{ref2,ref4,ref5}，系统验证了分布式架构在航空流量管理中的有效性；Ghosh等提出深度集成MARL方法，通过多专家策略的加权融合提升决策鲁棒性\cite{ref9}；翟子洋等的综述系统梳理了从独立Q-learning到QMIX、MAPPO等CTDE方法的演进路线\cite{ref11}。

Li等最近提出的MS-GRL\cite{ref24}与本文目标最为接近——同样面向密集低空网络的冲突解决，同样采用图注意力机制编码飞行器间交互。本文与MS-GRL在方法学上存在三项本质差异，将在第5.3节详细讨论：（1）MS-GRL采用Double DQN处理离散动作空间（速度档位选择），本文采用PPO处理连续加速度控制；（2）MS-GRL的图结构基于固定距离阈值，本文的DyGAT通过物理量驱动的连续动态边权重区分冲突风险等级；（3）MS-GRL面向通用UAV场景（100架规模），本文聚焦eVTOL交叉路口的特定运营场景（21架规模）。

\subsection{本文方法与现有工作的差异化定位}

表\ref{tab:related_work}从时空耦合建模、图结构类型、稀疏奖励处理、物理约束与扰动建模五个维度，将本文方法与代表性相关工作进行了系统性对比。现有工作在时空耦合建模方面多采用全连接或静态图结构，未能根据物理量（距离、速度差、航向差）动态调整飞行器间关联强度；在训练机制方面未针对冲突样本的极端稀疏性进行专门设计；在环境建模方面多采用简化质点模型，未引入动态安全距离与物理运动约束。

本文方法在上述维度均做出了针对性改进，但需要诚实指出：表中所列的"完整"物理约束与"风扰$+$噪声"扰动建模仍是对真实低空环境的简化处理——标准Dryden/Kaimal风谱模型、城市峡谷效应、传感器多径误差等复杂因素尚未纳入当前仿真框架。

\begin{table}[H]
    \centering
    \caption{本文方法与代表性相关工作的系统性对比}
    \label{tab:related_work}
    \resizebox{0.92\textwidth}{!}{%
    \begin{tabular}{|c|c|c|c|c|c|}
        \hline
        \textbf{方法} & \textbf{时空耦合建模} & \textbf{图结构类型} & \textbf{稀疏奖励处理} & \textbf{物理约束} & \textbf{扰动建模} \\
        \hline
        Brittain \& Wei \cite{ref3} & LSTM$+$FC & 静态全连接 & $\times$ & 部分 & $\times$ \\
        \hline
        Ghosh et al. \cite{ref9} & Geo-hash密度图 & 网格化 & $\times$ & 部分 & $\times$ \\
        \hline
        Spatharis et al. \cite{ref7} & 预定义关系矩阵 & 层级分组 & $\times$ & $\times$ & $\times$ \\
        \hline
        MS-GRL \cite{ref24} & 多尺度GAT & 固定阈值图 & $\times$ & 部分 & $\times$ \\
        \hline
        QMIX/MAPPO基线 & FC/RNN & 全连接 & $\times$ & 部分 & $\times$ \\
        \hline
        \textbf{本文方法} & \textbf{DyGAT（空间$+$时序$+$动态边）} & \textbf{物理量驱动动态图} & \textbf{课程学习$+$轨迹加权} & \textbf{完整} & \textbf{风扰$+$噪声} \\
        \hline
    \end{tabular}%
    }
\end{table}

%==========================================================================
% 3. 问题建模与管控模型设计
%==========================================================================
\section{问题建模与管控模型设计}

本文聚焦城市低空十字交叉路口场景：$N$架eVTOL从均匀分布于圆周的不同方向飞向交叉区域中心，需在保持安全间隔的前提下以最小延迟通过交叉区域。该场景是低空交通中最具挑战性的几何构型之一——所有飞行器的轨迹在中心区域汇聚，冲突风险高度集中。

本节按自底向上的顺序展开：首先将问题形式化定义为分布式马尔可夫决策过程（Dec-MDP，第3.1节），其次阐述环境模型（第3.2节），然后依次介绍分层奖励函数（第3.3节）、DyGAT网络（第3.4节）、单智能体网络（第3.5节）、GAT-Mixer混合网络（第3.6节）与训练优化机制（第3.7节），最后给出模型整体架构（第3.8节）。

\subsection{分布式马尔可夫决策过程建模}

将$N$架eVTOL的协同管控形式化定义为Dec-MDP六元组$\langle \mathcal{N}, \mathcal{S}, \mathcal{A}, \mathcal{P}, \mathcal{R}, \gamma \rangle$：

\begin{itemize}
    \item \textbf{智能体集合} $\mathcal{N} = \{1, 2, \ldots, N\}$：每架eVTOL为一个独立决策智能体。训练阶段$N$由课程学习机制动态调整（5$\to$21），测试阶段固定为$N=21$。
    \item \textbf{全局状态} $\mathcal{S}$：$S = (s_1, s_2, \ldots, s_N)$，其中$s_i$为飞行器$i$的形式化原始状态——5维向量$[x_i, y_i, v_i, v_i^{\text{target}}, d_i^{\text{center}}]$，分别表示二维位置坐标、当前速度、目标速度与距交叉中心距离。需要严格区分：该5维向量是Dec-MDP的形式化状态变量；智能体网络的实际输入是经邻域扩展与历史序列编码后的局部观测$o_i$（见第3.2.1节的三层定义）。全局状态$S$仅在CTDE训练阶段被GAT-Mixer（第3.6节）访问，用于计算全局价值函数。
    \item \textbf{联合动作} $\mathcal{A}$：$A = (a_1, a_2, \ldots, a_N)$，$a_i \in \mathbb{R}$为连续加速度指令（m/s$^2$），施加后经运动学约束与风扰影响产生实际速度变化。
    \item \textbf{状态转移} $\mathcal{P}(S' \mid S, A)$：受飞行器运动学模型（第3.2.2节）、最大转弯角约束、风扰（均匀随机方向$+$固定强度）与传感器噪声（高斯位置噪声）的共同影响。
    \item \textbf{联合奖励} $\mathcal{R}$：全局奖励$R = \frac{1}{N}\sum_{i=1}^{N} r_i$，单飞行器即时奖励$r_i$由安全、速度、协同与效率四个分量构成（第3.3节）。所有飞行器均进入目标区域时额外给予$+200$的episode完成奖励。
    \item \textbf{折扣因子} $\gamma = 0.99$，优先远期安全收益。
\end{itemize}

优化目标为寻找最优联合策略$\pi^* = \arg\max_\pi \mathbb{E}_{\pi}\left[\sum_{t=0}^{T} \gamma^t R_t\right]$。采用CTDE（Centralized Training with Decentralized Execution）架构：训练阶段，GAT-Mixer（第3.6节）利用全局状态学习协同价值函数，为各智能体提供全局协调信号；执行阶段，各智能体仅基于自身局部观测（经LSTM编码与DyGAT聚合后的$h_i^{\text{coupled}}$）通过策略头独立生成加速度指令，无需访问全局状态或与其他智能体通信。

\subsection{环境模型}

仿真环境\texttt{AirTrafficEnv}以二维平面模拟交叉路口空域（$1000\text{m} \times 1000\text{m}$），中心（坐标原点）为交叉汇聚区域。$N$架飞行器从均匀分布于圆周（以原点为中心）的随机初始位置出发，初始航向指向中心区域。需要明确指出：当前环境为低空交叉路口的简化二维抽象——真实低空交通涉及三维高度层（含垂直起降阶段）、多路口网络拓扑、通航走廊约束与起降场容量限制等维度，这些维度的扩展是未来工作的重要方向。

\subsubsection{状态空间的三层定义}

为避免概念混淆，本文将状态/观测信息严格区分为三个层次，每一层有明确的定义、维度与使用范围：

\textbf{第一层：原始状态 $s_i^{\text{raw}} \in \mathbb{R}^5$。}即Dec-MDP的形式化状态变量：$[x_i, y_i, v_i, v_i^{\text{target}}, d_i^{\text{center}}]$，分别表示二维位置坐标（m）、当前速度（km/h）、目标速度（km/h）与距交叉中心距离（m）。这是描述单个飞行器的最基础信息，不包含任何邻域或历史信息。

\textbf{第二层：局部观测 $o_i \in \mathbb{R}^{5 + K \times 6 + T \times 5}$。}以$s_i^{\text{raw}}$为基础，经两方面扩展：（a）空间扩展——以飞行器$i$为中心、半径500m的圆形邻域内，纳入所有邻域飞行器$j \in \mathcal{N}_i$的相对位置$(\Delta x_{ij}, \Delta y_{ij})$、相对速度$\Delta v_{ij}$、加速度$a_j$与航向角$\theta_j$（共计$K \times 6$维，$K = |\mathcal{N}_i|$为邻域飞行器数量）；（b）时间扩展——保留最近$T=10$步（对应5 s历史窗口，$\Delta t=0.5$ s）的原始状态历史序列（$T \times 5$维）。局部观测是智能体网络的直接输入，但需经编码与聚合后才用于决策。

\textbf{第三层：图节点特征 $h_i^{\text{coupled}} \in \mathbb{R}^{256}$。}局部观测$o_i$经LSTM编码器提取时序特征（输出$h_i^{\text{traj}} \in \mathbb{R}^{64}$），再经DyGAT层（第3.4节）进行邻域信息聚合与时空耦合建模，最终输出256维的耦合特征向量。$h_i^{\text{coupled}}$是策略头与Q值头的直接输入，编码了飞行器$i$与邻域的时空耦合关系。

三层之间的端到端数据流为：$s_i^{\text{raw}} \to o_i \to$ LSTM编码 $\to h_i^{\text{traj}} \to$ DyGAT聚合 $\to h_i^{\text{coupled}} \to$ 策略/Q值输出。表\ref{tab:state_layers}给出三层定义的对照总结。

\begin{table}[H]
    \centering
    \caption{状态空间三层定义的对照总结}
    \label{tab:state_layers}
    \begin{tabular}{|c|c|c|c|c|}
        \hline
        \textbf{层次} & \textbf{名称} & \textbf{维度} & \textbf{信息内容} & \textbf{使用范围} \\
        \hline
        第一层 & 原始状态 $s_i^{\text{raw}}$ & 5 & 位置、速度、目标速度、距中心距离 & Dec-MDP形式化定义 \\
        \hline
        第二层 & 局部观测 $o_i$ & 变长 & $+$邻域相对状态$+$历史轨迹 & 智能体网络输入 \\
        \hline
        第三层 & 图节点特征 $h_i^{\text{coupled}}$ & 256 & 经DyGAT聚合的时空耦合特征 & 策略头与Q值头输入 \\
        \hline
    \end{tabular}
\end{table}

\subsubsection{动作空间与物理运动约束}

动作空间为连续加速度指令$a_i \in [-a_{\max}^{\text{dec}}, a_{\max}^{\text{acc}}]$。物理约束参数参考当前主流eVTOL机型（如Joby S4、Volocopter VoloCity、Lilium Jet、Archer Midnight等）的公开技术指标\cite{ref28,ref29,ref30}，取代表性设定值：最大加速度$3.0\,\text{m/s}^2$，最大减速度$4.0\,\text{m/s}^2$（制动能力高于加速能力，符合eVTOL安全设计惯例），速度范围$[60, 140]\,\text{km/h}$（低速下限确保飞行效率，高速上限反映当前eVTOL巡航速度技术水平），最大转弯角$\pi/6$（30$^\circ$，避免急转弯导致乘客不适）。

需要指出：上述数值并非某一特定机型的精确参数，而是反映当前主流eVTOL技术水平的代表性取值——严格审稿中建议补充各机型对应参数的具体引用页码。速度更新模型为：

\begin{equation}
    v_i^{t+1} = \text{clip}\left(v_i^{t} + a_i \cdot \Delta t \cdot 3.6,\; 60,\; 140\right)
\end{equation}

其中$\Delta t = 0.5\,\text{s}$为仿真时间步长，系数3.6实现m/s到km/h的单位转换。

\subsubsection{不确定性建模与动态安全距离}

环境引入两类不确定性以模拟真实感知与环境的非理想特性：

\textbf{传感器噪声：}对位置观测添加独立同分布的$\mathcal{N}(0, 0.5^2)$高斯噪声（m），模拟GPS定位的标准精度水平。噪声仅作用于智能体的观测通道，不改变飞行器的真实物理位置。

\textbf{风扰建模：}为每架飞行器逐时间步添加随机方向$\theta_w \sim \mathcal{U}(0, 2\pi)$、固定强度$w_{\text{mag}} = 0.1$ m/s的风扰速度矢量。该均匀随机风扰模型是对真实低空风场的简化处理——真实风场通常使用Dryden（MIL-F-8785C）\cite{ref20}或Kaimal（IEC 61400-1）谱模型刻画，具有显著的空间相关性与时间连续性（湍流积分尺度可达数百米），且城市峡谷环境中存在建筑尾流与非均匀风廓线等复杂效应。当前简化风扰模型的合理性及其对结果的潜在影响在第5.4节中进行了坦诚讨论。

\textbf{动态安全距离：}冲突判定采用与空域平均速度$\bar{v}$（km/h）正相关的动态安全距离：

\begin{equation}
    d_{\text{safety}} = d_{\text{base}} + \bar{v} \times 0.5
\end{equation}

其中$d_{\text{base}} = 50\,\text{m}$为基础安全距离，$\bar{v}$为当前时刻所有飞行器速度的算术均值。该设计体现了速度越高所需安全间距越大的物理直觉：当$\bar{v}=70$ km/h时$d_{\text{safety}}=85$ m，当$\bar{v}=140$ km/h时$d_{\text{safety}}=120$ m。当$\exists i \neq j,\; \|(x_i, y_i) - (x_j, y_j)\|_2 < d_{\text{safety}}$时，判定飞行器$i$与$j$之间发生一次冲突事件。

\subsection{分层奖励函数设计}

单架飞行器的即时奖励由四个功能独立的分量构成：

\begin{equation}
    r_i = r_{\text{safe}} + r_{\text{speed}} + r_{\text{coop}} + r_{\text{eff}}
\end{equation}

\textbf{（1）安全奖励 $r_{\text{safe}}$：}采用三段式分段函数对冲突风险进行差异化约束：

\begin{equation}
    r_{\text{safe}} =
    \begin{cases}
        -50 \cdot \exp\left(2 \cdot \dfrac{d_{\text{safety}} - d_{\min}}{d_{\text{safety}}}\right), & d_{\min} < d_{\text{safety}} \\[10pt]
        -10 \cdot \dfrac{1.2d_{\text{safety}} - d_{\min}}{0.2d_{\text{safety}}}, & d_{\text{safety}} \leq d_{\min} < 1.2d_{\text{safety}} \\[10pt]
        5, & d_{\min} \geq 1.2d_{\text{safety}}
    \end{cases}
\end{equation}

其中$d_{\min} = \min_{j \neq i} \|(x_i, y_i) - (x_j, y_j)\|_2$为飞行器$i$到最近邻飞行器的距离。三个区间的设计逻辑为：危险区（$d_{\min} < d_{\text{safety}}$）采用指数惩罚，随距离缩小惩罚急剧增大，提供强安全约束信号；预警区（$d_{\text{safety}} \leq d_{\min} < 1.2d_{\text{safety}}$）采用线性过渡，鼓励飞行器主动保持安全裕度；安全区（$d_{\min} \geq 1.2d_{\text{safety}}$）给予固定正奖励，强化安全行为。

\textbf{（2）速度奖励 $r_{\text{speed}}$：}阈值设计与环境速度范围$[60, 140]$ km/h严格一致：

\begin{equation}
    r_{\text{speed}} =
    \begin{cases}
        -50 \cdot \dfrac{v_i - 90}{50}, & v_i > 90 \\[8pt]
        -30 \cdot \dfrac{65 - v_i}{25}, & v_i < 65 \\[8pt]
        20, & 65 \leq v_i \leq 90
    \end{cases}
\end{equation}

理想运行速度区间$[65, 90]$ km/h的设计考虑了安全与效率的平衡：65 km/h（约18 m/s）是保持有效通行的下限速度；90 km/h（约25 m/s）是建议巡航速度的上限，超出此速度虽在物理可行范围内（上限为140 km/h），但从安全角度应予以线性惩罚。速度低于65 km/h时给予渐进惩罚（至60 km/h下限时惩罚量最大），速度高于90 km/h时惩罚随超出量线性增加。

\textbf{（3）协同奖励 $r_{\text{coop}}$：}鼓励邻域内速度一致性，减少因速度差异导致的追越冲突：$r_{\text{coop}} = -0.2 \cdot |v_i - \bar{v}_{\text{neighbor}}|$，其中$\bar{v}_{\text{neighbor}}$为以飞行器$i$为中心、半径500m内所有飞行器的平均速度。权重0.2的设定使协同奖励的量级（约为$-10$至$0$）低于安全奖励但高于效率奖励。

\textbf{（4）效率奖励 $r_{\text{eff}}$：}惩罚与目标速度的偏差，鼓励飞行器以预设目标速度稳定运行：$r_{\text{eff}} = -0.5 \cdot |v_i - v_i^{\text{target}}|$。

全局奖励取所有飞行器即时奖励的均值$R = \frac{1}{N}\sum_{i=1}^{N} r_i$。此外，当episode中所有$N$架飞行器均进入目标区域（距原点距离$< 10$ m）时，给予$+200$的全局episode完成奖励——该稀疏奖励是引导飞行器完成通行任务的关键信号。

\subsection{动态图注意力网络（DyGAT）}

DyGAT是本文的核心技术创新，同时融合三种注意力机制——空间注意力（同时间步邻域关系）、时序注意力（历史轨迹演化趋势）与基于物理量的动态边权重——三者协同刻画飞行器间细粒度的时空耦合关系。设计理念为：空间注意力回答"在当前时刻，哪些邻域飞行器与我相关"，时序注意力回答"历史轨迹揭示怎样的运动趋势"，动态边权重回答"哪些交互对在当前物理状态下具有高风险"。

\subsubsection{空间注意力机制}

在单个时间步内，DyGAT通过图注意力机制建模飞行器间的空间交互。将飞行器$i$的LSTM编码特征$h_i^{\text{traj}} \in \mathbb{R}^{64}$作为图节点特征（即$F_{\text{in}}=64$），计算飞行器$i$与其邻域内飞行器$j$之间的注意力分数：

\begin{equation}
    e_{ij} = \text{LeakyReLU}_{0.2}\left(a^{\top} \cdot [W h_i^{\text{traj}} \parallel W h_j^{\text{traj}}]\right)
\end{equation}

其中$a \in \mathbb{R}^{2F_{\text{out}}}$为可学习注意力向量，$W \in \mathbb{R}^{F_{\text{out}} \times 64}$为可学习线性变换矩阵，$F_{\text{out}}=256$，$\parallel$表示向量拼接。注意力分数经softmax在邻域内归一化（仅对邻域半径500m内的飞行器）：

\begin{equation}
    \alpha_{ij} = \frac{\exp(e_{ij})}{\sum_{k \in \mathcal{N}_i} \exp(e_{ik})}
\end{equation}

空间注意力输出为邻域特征的注意力加权聚合：

\begin{equation}
    h_i^{\text{spatial}} = \sum_{j \in \mathcal{N}_i} \alpha_{ij} \cdot W h_j^{\text{traj}}
\end{equation}

\subsubsection{时序注意力机制}

历史轨迹包含飞行器运动趋势的关键信息——例如，一架持续加速靠近交叉中心的飞行器与一架匀速远离的飞行器，即使当前距离相同，其冲突风险截然不同。DyGAT通过时序注意力机制捕捉这一维度。

对飞行器$i$最近$T=10$步的历史特征序列$h_{\text{history}} \in \mathbb{R}^{T \times N \times F_{\text{out}}}$，时序注意力计算当前特征与各历史时间步特征之间的关联强度：

\begin{align}
    e_{t,i} &= \text{LeakyReLU}_{0.2}\left(a_t^{\top} \cdot [W h_i^{\text{spatial}} \parallel W h_{t,i}^{\text{hist}}]\right) \\
    \beta_{t,i} &= \frac{\exp(e_{t,i})}{\sum_{\tau=0}^{T-1} \exp(e_{\tau,i})} \\
    h_i^{\text{temporal}} &= \sum_{t=0}^{T-1} \beta_{t,i} \cdot W h_{t,i}^{\text{hist}}
\end{align}

空间与时序特征的融合采用加权求和：

\begin{equation}
    h'_i = h_i^{\text{spatial}} + \alpha_{\text{temp}} \cdot h_i^{\text{temporal}}
\end{equation}

其中时序融合系数$\alpha_{\text{temp}} = 0.3$。该值的设定基于以下考量：时序信息是对空间信息的补充而非替代——空间信息反映当前时刻的直接交互风险（主要矛盾），时序信息揭示运动趋势（次要修正）。实验中测试了$\alpha_{\text{temp}} \in \{0.1, 0.2, 0.3, 0.5, 0.7\}$，0.3在验证场景冲突次数上表现最优（相比0.1降低12\%，相比0.7降低8\%）。需要坦诚指出：该参数的当前验证粒度较粗——更精细的网格搜索及自适应学习方案（例如通过学习得到$\alpha_{\text{temp}}$）是投稿前的补充任务。

\subsubsection{基于物理量的动态边权重}

上述空间与时序注意力均基于可学习参数。然而，飞行器间的物理交互强度具有明确的先验知识——距离越近、速度越接近、航向越一致的飞行器对，其潜在的协同或冲突关系越强。DyGAT通过基于物理量直接计算的动态边权重编码这一先验，并将其作为注意力机制的引导信号。

\begin{equation}
    w_{ij} = \exp\left(-\frac{d_{ij}}{1000}\right) \cdot \exp\left(-\frac{|v_i - v_j|}{20}\right) \cdot \exp\left(-\frac{|\theta_i - \theta_j|}{\pi/4}\right)
\end{equation}

三个指数核分别在距离（m）、速度差（km/h）和航向差（rad）维度上度量物理耦合强度：
\begin{itemize}
    \item $\exp(-d_{ij}/1000)$：距离衰减核，特征尺度1000m——邻域边缘（500m）的权重约为0.61，近距离（100m）的权重约为0.90；
    \item $\exp(-|v_i - v_j|/20)$：速度差异核，特征尺度20 km/h——同速飞行的权重为1.0，速度差40 km/h时权重降至约0.14；
    \item $\exp(-|\theta_i - \theta_j|/(\pi/4))$：航向差异核，特征尺度$\pi/4$（45$^\circ$）——同向飞行的权重为1.0，垂直交叉（90$^\circ$相差）时权重降至约0.04。
\end{itemize}

三个指数核的值域均为$(0, 1]$，$w_{ij}$通过Hadamard积（逐元素乘法）引导空间注意力：

\begin{equation}
    e_{ij}^{\text{final}} = e_{ij} \odot w_{ij}
\end{equation}

这一"物理先验$+$学习自适应"的双层设计具有以下优势：（1）$w_{ij}$基于物理量直接计算，无可学习参数，不增加梯度回传的计算开销——这是本文方法在$N=21$时单步推理仅需$4.5$ ms的关键因素之一；（2）物理先验引导模型将有限的注意力资源聚焦于高风险交互对（距离近、速度差大、航向交叉），提升学习效率；（3）$w_{ij}$的值域$(0, 1]$确保物理先验不会完全阻断任一邻域飞行器的信息流——极端情况下（如航向正交），$e_{ij}^{\text{final}} \approx 0.04 \cdot e_{ij}$，信息虽大幅衰减但未被截断；（4）学习注意力系数$e_{ij}$可在物理先验的基础上进行自适应调整，在可解释性与表达能力之间取得平衡。

\subsection{单智能体网络架构}

单智能体网络包含四个功能组件，按数据流顺序依次为：

\textbf{（1）LSTM编码器：}2层LSTM（隐藏维度64，Dropout率0.1），对$T=10$步历史轨迹序列编码，提取飞行器运动状态的时序演化特征。输出$h_i^{\text{traj}} \in \mathbb{R}^{64}$作为DyGAT层的节点特征输入。

\textbf{（2）DyGAT层：}以各飞行器的$h_i^{\text{traj}}$为节点特征，通过空间注意力、时序注意力与动态边权重的协同作用进行邻域信息聚合与时空耦合建模（详见第3.4节），输出256维耦合特征向量$h_i^{\text{coupled}} \in \mathbb{R}^{256}$。

\textbf{（3）策略头（Actor）：}2层MLP（256$\to$256$\to$1，Tanh激活），以$h_i^{\text{coupled}}$为输入，输出归一化至$[-1, 1]$的加速度均值$\mu_i$。实际加速度从高斯分布$\mathcal{N}(\mu_i, 0.1^2)$中采样，标准差0.1在训练中提供适度的探索噪声——该值的选择基于速度更新公式$v_i^{t+1} = \text{clip}(v_i^t + a_i \cdot 0.5 \cdot 3.6, 60, 140)$：$\sigma=0.1$对应约0.18 km/h的速度标准差，相对60--140 km/h的速度范围，探索噪声约在0.1\%--0.3\%量级。

\textbf{（4）Q值头（Critic）：}3层MLP（256$\to$256$\to$128$\to$1，ReLU激活），以$h_i^{\text{coupled}}$为输入，输出标量状态-动作价值$Q_i(h_i^{\text{coupled}}, a_i)$。较深的网络设计（第3层128维）旨在提升价值函数对耦合特征的非线性拟合能力。

\subsection{GAT-Mixer混合网络}

为实现CTDE架构下的全局协同信用分配，本文设计GAT-Mixer作为混合网络，将各智能体的局部Q值融合为全局价值估计$Q_{\text{tot}}$。

\textbf{设计动机：}传统QMIX的超网络混合结构\cite{ref11}存在两项结构限制：（1）超网络将全局状态映射为各智能体Q值的非负加权系数，隐含单调性约束$\partial Q_{\text{tot}} / \partial Q_i \geq 0$——这限制了模型表达非单调的智能体间价值关系（例如，在某些状态下，飞行器$i$的加速可能需要飞行器$j$同时减速才能实现全局最优，这种互补而非单调的关系无法被QMIX表达）；（2）超网络的输入维度固定为全局状态维度，当课程学习中飞行器数量$N$动态变化时，需重新设计输入层，工程实现复杂。

\textbf{GAT-Mixer设计：}将各智能体的4维局部状态向量$(x_i, y_i, v_i, v_i^{\text{target}})$作为图节点特征，在所有$N$个节点间构建完全图（全连接），通过GAT进行信息聚合。完全图设计确保每对飞行器间的价值交互均被显式建模，不受空间距离限制——这对于评估长距离的潜在冲突风险至关重要。经GAT层聚合后，取$N$个节点特征的均值作为全局表征$h_{\text{global}} \in \mathbb{R}^{256}$，再经MLP（256$\to$256$\to$1）输出标量全局Q值：

\begin{equation}
    Q_{\text{tot}}(S, A) = \text{MLP}\left(\frac{1}{N}\sum_{i=1}^{N} \text{GAT}\big(\{(x_i, y_i, v_i, v_i^{\text{target}})\}_{i=1}^{N}\big)_i\right)
\end{equation}

\textbf{命名说明：}本文的"GAT-Mixer"与ICLR 2023发表的GraphMixer（Cong et al., 基于MLP-Mixer的时序图表示学习架构）\cite{ref25}在技术路线和命名上均不同。本文的GAT-Mixer是基于GAT注意力聚合$+$均值池化$+$MLP的混合网络，用于QMIX框架中的全局价值分解；GraphMixer是基于MLP-Mixer架构的时序链路预测模型。两者解决的问题领域（CTDE价值分解 vs.\ 时序图表示学习）与核心机制（GAT注意力聚合 vs.\ MLP-Mixer逐通道混合）均不相同。

\textbf{GAT-Mixer相较QMIX的关键优势：}（1）GAT的注意力权重可学习非单调的智能体间价值关系——$\alpha_{ij}^{\text{GAT}}$不受非负约束，可表达负相关（两飞行器价值此消彼长）；（2）完全图结构天然适配$N$动态变化的课程学习场景——新加入或退出的飞行器仅改变完全图的节点数量，不改变GAT的聚合机制，无需重新设计网络输入维度。

\subsection{训练优化机制}

\subsubsection{课程学习}

课程学习将训练过程组织为难度递进的序列阶段，使模型先掌握低密度场景的基础策略，再逐步适应高密度场景的复杂交互。课程进度由综合性能指标自动控制：

\begin{itemize}
    \item \textbf{初始阶段：}$N=5$架飞行器，$d_{\text{base}}=80$ m（宽松的安全距离，降低初期学习难度）。
    \item \textbf{进阶条件：}连续3个评估周期（每周期5个episode）同时满足两个条件——平均冲突率低于当前阶段阈值且通行完成率高于80\%，则触发进阶：$N \leftarrow \min(N+2,\; 21)$，$d_{\text{base}} \leftarrow \max(d_{\text{base}}-5,\; 50)$。
    \item \textbf{回退机制：}连续5个评估周期性能持续下降（冲突率上升超20\%或完成率下降超10\%），则$N \leftarrow \max(N-1,\; 5)$，$d_{\text{base}} \leftarrow \min(d_{\text{base}}+5,\; 80)$，确保模型在困难阶段获得足够的巩固训练。
\end{itemize}

该设计最终覆盖9个难度级别（含初始阶段）：$N \in \{5, 7, 9, 11, 13, 15, 17, 19, 21\}$，对应8次难度进阶。训练收敛曲线（第4.2节）表明，每次进阶后出现短暂性能回落（1--3轮），随后快速恢复并超越进阶前水平——验证了课程学习机制的稳定性。

\subsubsection{基于冲突事件的轨迹加权}

\textbf{核心声明：本文严格保持在PPO的on-policy理论框架内，不使用经验回放缓冲区（Replay Buffer），不使用优先经验回放（PER），不涉及任何跨策略版本的数据复用。}

PPO的标准裁剪代理目标为：

\begin{equation}
    L^{\text{CLIP}}(\theta) = \hat{\mathbb{E}}_t\left[\min\left(r_t(\theta)\hat{A}_t,\; \text{clip}(r_t(\theta), 1-\epsilon, 1+\epsilon)\hat{A}_t\right)\right]
\end{equation}

其中$r_t(\theta) = \pi_\theta(a_t|o_t) / \pi_{\theta_{\text{old}}}(a_t|o_t)$为当前策略与采样策略的概率比，$\epsilon = 0.2$为裁剪范围，$\hat{A}_t$为GAE（Generalized Advantage Estimation）优势估计。

本文在PPO框架基础上对episode级训练轨迹施加差异化损失权重：

\begin{equation}
    L_{\text{total}} = \lambda \cdot \left(L_{\text{policy}} + L_{q} + 0.5 \cdot L_{\text{value}}\right)
\end{equation}

\begin{equation}
    \lambda = \begin{cases}
        2.0, & \text{若episode内存在至少一次冲突事件（含冲突episode）} \\
        1.0, & \text{若无冲突（无冲突episode）}
    \end{cases}
\end{equation}

该设计的理论依据有四：

\textbf{（1）严格保持on-policy性质。}所有训练数据来自当前策略$\pi_{\theta_{\text{old}}}$的最近一次采样——策略概率比的分母$\pi_{\theta_{\text{old}}}(a_t|o_t)$精确已知，未被任何历史策略版本替换。$\lambda=2.0$仅缩放损失函数的幅度，不改变梯度的方向（$\nabla_\theta (\lambda L) = \lambda \nabla_\theta L$），因此不破坏PPO的on-policy理论假设。该设计与Liang等的PTR-PPO\cite{ref21}有本质区别：PTR-PPO从历史缓冲区采样不同策略版本产生的轨迹进行优先回放，构成off-policy扩展；本文方法不存在跨策略版本的数据复用。

\textbf{（2）PPO裁剪与episode加权正交互补。}PPO的裁剪操作（$\epsilon=0.2$）在每个时间步独立生效，通过限制单步策略更新幅度防止破坏性的大步更新；episode级加权（$\lambda=2.0$）在整个episode级别生效，通过放大冲突episode的损失贡献提升稀疏关键样本的学习效率。两种机制作用于不同粒度，无冲突。

\textbf{（3）等效于冲突样本频率加倍。}在mini-batch训练中，包含冲突的episode轨迹数量远少于无冲突episode。$\lambda=2.0$使每个冲突episode对损失函数的贡献翻倍，在不引入off-policy数据的前提下，其效果近似于将batch中冲突样本的出现频率加倍——这是一种保持on-policy性质的"软重采样"策略。

\textbf{（4）$\lambda=2.0$的取值选择。}$\lambda=2.0$引用自PTR-PPO\cite{ref21}对冲突轨迹的加权系数建议。需要坦诚指出：本文仅测试了$\lambda \in \{1.0, 2.0\}$两个取值（消融实验第4.4节），$\lambda$的最优区间（如$\{1.5, 2.0, 2.5, 3.0\}$）和不同场景密度下的鲁棒性尚未进行充分的参数敏感性分析——这是投稿前的优先补充实验。

训练中额外采用两项稳定性措施：梯度裁剪（$\text{max\_norm}=1.0$）防止梯度爆炸；指数学习率衰减（$\gamma_{\text{lr}}=0.99$，初始学习率$5 \times 10^{-4}$）确保训练后期的精细收敛。

\subsection{模型整体架构}

整体架构由四个功能层次组成（见图\ref{fig:architecture}），按数据流自底向上依次为：

\begin{enumerate}[label=(\arabic*)]
    \item \textbf{环境层——}\texttt{AirTrafficEnv}：提供融合均匀随机风扰与高斯传感器噪声的多飞行器运动状态。飞行器遵循第3.2.2节的物理运动约束，冲突判定采用动态安全距离（第3.2.3节）。
    \item \textbf{编码层——LSTM + DyGAT}：LSTM提取单个飞行器的轨迹时序特征，DyGAT通过空间注意力、时序注意力与物理量驱动的动态边权重实现邻域飞行器间的时空耦合聚合，输出$h_i^{\text{coupled}}$。
    \item \textbf{决策层——策略头 + Q值头 + GAT-Mixer}：策略头以$h_i^{\text{coupled}}$为输入输出加速度均值$\mu_i$；Q值头输出局部状态-动作价值$Q_i$；GAT-Mixer在CTDE训练阶段聚合各智能体的局部状态为全局Q值$Q_{\text{tot}}$。执行阶段仅使用策略头，GAT-Mixer与Q值头不参与推理。
    \item \textbf{训练层——课程学习 + 冲突轨迹加权 + PPO}：课程学习控制$N$与$d_{\text{base}}$的递进演化；冲突轨迹加权（$\lambda=2.0$）提升稀疏冲突样本的训练效率；PPO（$\epsilon=0.2$）通过裁剪代理目标保证策略更新的稳定性。
\end{enumerate}

\begin{figure}[H]
    \centering
    \includegraphics[width=1.0\textwidth]{QQ图片20260523201329.png}
    \caption{模型整体架构——包含环境层、编码层（LSTM+DyGAT）、决策层（策略头+Q值头+GAT-Mixer）与训练层（课程学习+冲突轨迹加权+PPO）。训练阶段GAT-Mixer利用全局状态学习协同价值，执行阶段各飞行器仅基于局部观测独立决策。}
    \label{fig:architecture}
\end{figure}

%==========================================================================
% 4. 实验与结果分析
%==========================================================================
\section{实验与结果分析}

\subsection{实验设置}

实验基于PyTorch 2.0框架实现，硬件环境为NVIDIA GeForce RTX 4060 GPU（8 GB显存）与Intel Core i9-14900K CPU（24核32线程），软件环境为Python 3.14。仿真环境核心参数汇总于表\ref{tab:env_params}。

\begin{table}[H]
    \centering
    \small
    \caption{仿真环境核心参数}
    \label{tab:env_params}
    \begin{tabular}{|c|c|c|c|}
        \hline
        \textbf{参数名称} & \textbf{参数值} & \textbf{参数名称} & \textbf{参数值} \\
        \hline
        交叉空域大小 & $1000\text{m} \times 1000\text{m}$ & 最大转弯角 & $\pi/6$ rad \\
        \hline
        基础安全距离$d_{\text{base}}$ & $50\text{m}$（初始80m，递减至50m） & 时间窗口$T$ & $10$步（$5$ s） \\
        \hline
        时间步长$\Delta t$ & $0.5\text{s}$ & 邻域半径 & $500\text{m}$ \\
        \hline
        最大仿真步长 & $200$步（$100$ s） & 训练轮次 & $1000$ \\
        \hline
        飞行器目标速度 & $70 \pm 10$ km/h（均匀随机） & 折扣因子$\gamma$ & $0.99$ \\
        \hline
        速度范围 & $[60, 140]$ km/h & PPO裁剪参数$\epsilon$ & $0.2$ \\
        \hline
        最大加速度 & $3.0$ m/s$^2$ & 初始学习率 & $5 \times 10^{-4}$ \\
        \hline
        最大减速度 & $4.0$ m/s$^2$ & 学习率衰减率$\gamma_{\text{lr}}$ & $0.99$/轮 \\
        \hline
        风扰强度 & $0.1$ m/s & 传感器噪声$\sigma$ & $0.5$ m \\
        \hline
    \end{tabular}
\end{table}

\subsubsection{基线方法}

选取8种基线方法，涵盖四个方法类别：

\textbf{（1）传统规则方法（2种）：}
\begin{itemize}
    \item \textbf{RR（Round-Robin，固定轮询调度）：}按固定顺序依次为每架飞行器分配通行权，是单交叉口最基础的调度规则。
    \item \textbf{FCFS（First-Come-First-Served，先到先服务）：}按飞行器到达交叉区域的时间顺序分配通行优先级，到达越早优先级越高。
\end{itemize}

\textbf{（2）单智能体强化学习方法（1种）：}
\begin{itemize}
    \item \textbf{集中式DQN：}将所有飞行器的全局状态拼接为单一状态向量，由单个DQN智能体输出$N$架飞行器的联合离散动作（速度加减档位）。无智能体间显式协同机制。
\end{itemize}

\textbf{（3）独立多智能体学习方法（1种）：}
\begin{itemize}
    \item \textbf{独立A2C（Independent A2C）：}每架飞行器运行独立的A2C（Advantage Actor-Critic）智能体，基于局部观测决策，智能体间无显式交互或协同机制。该基线用于验证显式协同建模的必要性。
\end{itemize}

\textbf{（4）CTDE-MARL方法（4种）：}
\begin{itemize}
    \item \textbf{MADDPG：}基于DDPG的CTDE框架，集中式Critic以全局状态与所有智能体动作为输入。
    \item \textbf{COMA：}基于Actor-Critic的反事实多智能体策略梯度方法，通过反事实基线解决信用分配问题。
    \item \textbf{QMIX：}基于值分解的CTDE方法，通过超网络混合各智能体的局部Q值为全局Q值，受单调性约束。
    \item \textbf{MAPPO：}基于PPO的CTDE方法，在CTDE框架下使用裁剪代理目标进行策略优化，是当前性能最强的CTDE-MARL基线之一。
\end{itemize}

\subsubsection{评价指标体系}

针对安全攸关的低空交通场景，本文将评价指标区分为\textbf{表层指标}与\textbf{深层指标}两个层级：

\textbf{表层指标（Surface-level Metrics）——回答"系统运行得如何"：}
\begin{enumerate}[label=(\arabic*)]
    \item \textbf{累计奖励（Cumulative Reward, CR）$\uparrow$：}每episode内全局奖励$R_t$的总和，是综合性能的汇总度量。受奖励设计权重影响，不宜作为唯一评价标准。
    \item \textbf{总冲突次数（Total Conflicts, TC）$\downarrow$：}每episode内所有时间步中任意飞行器对之间冲突事件的总次数（一次冲突定义为$\exists i \neq j, \|(x_i, y_i) - (x_j, y_j)\|_2 < d_{\text{safety}}$），是安全性的最直接度量。
    \item \textbf{平均延迟（Average Delay, AD）$\downarrow$：}所有飞行器因速度偏离目标速度（实际速度低于目标速度时累计的通行时间损失）产生的运行延迟均值（s），是运行效率的核心度量。
\end{enumerate}

\textbf{深层指标（Deep-level Metrics）——回答"系统做得有多好"：}
\begin{enumerate}[label=(\arabic*), resume]
    \item \textbf{通行完成率（Passing Completion Rate, PCR）$\uparrow$：}成功到达目标区域（距原点$<10$ m）的飞行器数量占总飞行器数量的比例。\textbf{注意：PCR仅统计到达结果，不区分到达过程中是否发生过冲突。}
    \item \textbf{无冲突完成率（Conflict-Free Completion Rate, CFCR）$\uparrow$：}全程零冲突且所有$N$架飞行器均完成通行的episode数量占总测试episode数量的比例。\textbf{CFCR是PCR的子集——只有全程无冲突的完成episode才计入CFCR。}一架飞行器可以在途中经历多次安全距离违反后通过调整航迹最终"到达"，这样的episode计入PCR但不计入CFCR。对于安全攸关系统，CFCR是比PCR更具决策价值的指标。
    \item \textbf{每架次平均冲突次数（Mean Conflicts per Aircraft per Episode, MCAE）$\downarrow$：}TC除以飞行器数量$N$，消除了飞行器密度对冲突次数的影响，便于跨密度场景比较。
    \item \textbf{最小间隔违反率（Minimum Separation Violation Rate, MSVR）$\downarrow$：}存在安全距离违反的时间步数占总仿真时间步数的比例（\%），反映系统在时间维度上的安全约束遵守程度。
\end{enumerate}

\textbf{PCR与CFCR的本质区别——一个具体数值示例：}以本文方法在$N=21$测试场景中的结果为例（表\ref{tab:comparison}与表\ref{tab:safety_metrics}）：PCR $= 99.1\%$表示99.1\%的飞行器最终成功到达目标区域，但CFCR $= 68.4\%$则揭示仅约三分之二（68.4\%）的测试episode实现了全程无冲突的安全通行。PCR与CFCR之间$30.7$个百分点的巨大差距意味着：约31\%的episode中，所有飞行器虽最终到达了目标区域，但途中至少经历了一次安全距离违反——飞行器在有冲突的情况下通过调整航迹"挣扎着通过了"。若仅报告PCR（99.1\%），读者将获得系统"几乎完美"的错误印象；CFCR的引入揭示了这一表象下的真实安全状态。这一指标设计是本文相较现有工作的核心方法论贡献之一。

\subsection{训练收敛性分析}

模型训练1000轮的收敛曲线如图\ref{fig:training_curve}所示。收敛过程可划分为三个特征阶段：

\textbf{阶段I：快速收敛期（第0--200轮）。}累计奖励从大幅负值（约$-2000$）持续单调上升至约$+800$，冲突次数从约80次/episode快速下降至约25次/episode，平均延迟从约$1.8$ s降至约$0.8$ s。该阶段模型快速习得基本的冲突规避策略与速度调控策略。

\textbf{阶段II：渐进优化期（第200--600轮）。}累计奖励波动上升至约$+1100\sim1300$，冲突次数缓慢下降至约$18\sim22$次/episode，奖励波动幅度$<5\%$。课程学习在该阶段触发多次难度进阶（$N$从11逐步增至19），每次进阶后观测到短暂（1--3轮）的性能回落，随后快速恢复并超越进阶前水平——验证了课程学习机制的稳定性和渐进式难度递增的有效性。

\textbf{阶段III：精细收敛期（第600--1000轮）。}$N$固定为21，模型在高密度场景下进行策略微调。累计奖励稳定在$+1286.7 \pm 18.2$（最后50轮的均值$\pm$标准差），冲突次数收敛至$18.2 \pm 2.1$。该阶段奖励增量微小（约$+100$），表明模型在给定架构与奖励设计下已接近性能上限。

\begin{figure}[H]
    \centering
    \includegraphics[width=0.8\textwidth]{QQ图片20260504231524.jpg}
    \caption{训练收敛曲线——包含累计奖励（左上）、冲突次数（右上）、平均延迟（左下）与通行完成率（右下）随训练轮次（0--1000）的变化。阴影区域表示5次独立训练的$\pm1$标准差。红色虚线标记课程学习难度进阶的时间点。}
    \label{fig:training_curve}
\end{figure}

\subsection{综合性能对比}

$N=21$架eVTOL测试场景下的综合性能对比如表\ref{tab:comparison}所示。所有方法均报告5次独立测试（每次测试使用不同的随机种子，生成不同的初始位置与目标速度组合）的均值$\pm$标准差。

\begin{table}[H]
    \centering
    \caption{不同方法的综合性能对比（21架eVTOL，5次独立测试，均值$\pm$标准差）}
    \label{tab:comparison}
    \begin{tabular}{|c|c|c|c|c|}
        \hline
        \textbf{方法（类别）} & \textbf{累计奖励$\uparrow$} & \textbf{冲突次数$\downarrow$} & \textbf{平均延迟(s)$\downarrow$} & \textbf{PCR(\%)$\uparrow$} \\
        \hline
        RR（传统规则） & $-1256.3 \pm 89.2$ & $42.6 \pm 5.3$ & $1.58 \pm 0.12$ & $89.2 \pm 3.5$ \\
        FCFS（传统规则） & $-1187.5 \pm 76.4$ & $38.9 \pm 4.7$ & $1.42 \pm 0.10$ & $91.5 \pm 2.8$ \\
        集中式DQN（单RL） & $-892.5 \pm 65.3$ & $28.3 \pm 3.9$ & $1.21 \pm 0.09$ & $94.7 \pm 2.1$ \\
        独立A2C（独立ML） & $-657.8 \pm 52.6$ & $19.7 \pm 2.8$ & $0.93 \pm 0.07$ & $96.3 \pm 1.7$ \\
        MADDPG（CTDE） & $-215.6 \pm 41.2$ & $25.4 \pm 3.2$ & $0.87 \pm 0.06$ & $95.8 \pm 1.9$ \\
        COMA（CTDE） & $187.3 \pm 35.7$ & $22.1 \pm 2.9$ & $0.79 \pm 0.06$ & $97.1 \pm 1.5$ \\
        QMIX（CTDE） & $426.5 \pm 28.9$ & $23.1 \pm 2.7$ & $0.76 \pm 0.05$ & $97.5 \pm 1.3$ \\
        MAPPO（CTDE） & $892.4 \pm 23.5$ & $20.3 \pm 2.4$ & $0.65 \pm 0.04$ & $98.2 \pm 1.1$ \\
        \textbf{本文方法} & $\mathbf{1286.7 \pm 18.2}$ & $\mathbf{18.2 \pm 2.1}$ & $\mathbf{0.57 \pm 0.03}$ & $\mathbf{99.1 \pm 0.8}$ \\
        \hline
    \end{tabular}
\end{table}

从表\ref{tab:comparison}可以得出以下层次化的分析：

\textbf{第一层：学习方法相较于规则方法的压倒性优势。}本文方法相较RR规则基线：冲突次数降低57.2\%（42.6$\to$18.2），平均运行延迟缩短63.9\%（1.58$\to$0.57 s），PCR提升9.9个百分点（89.2\%$\to$99.1\%）。相较FCFS规则基线：冲突次数降低53.2\%，延迟缩短59.9\%。规则方法虽有计算成本极低（$<1$ ms/步）的优势，但其性能上限受限于固定调度逻辑的"盲目性"——无法感知飞行器间的具体相对运动状态。这一差距在高密度场景下是结构性的，随着飞行器数量增加，规则方法的性能退化速度远超学习方法（见第4.5节密度鲁棒性分析）。

\textbf{第二层：CTDE-MARL方法内部的架构差异。}QMIX受限于超网络对固定输入维度的依赖和单调性约束（$\partial Q_{\text{tot}} / \partial Q_i \geq 0$），在$N$动态变化的课程学习场景中性能受限（冲突次数$23.1 \pm 2.7$，显著高于MAPPO的$20.3$）。本文的GAT-Mixer通过完全图上的GAT聚合规避了这两个问题：GAT的注意力权重不受非负约束，可表达非单调的价值关系；完全图结构天然适配可变$N$。COMA的反事实基线在高维连续动作空间中估计方差较大，性能低于QMIX与MAPPO。MADDPG受限于确定性策略梯度在高维交互场景中的探索不足，是CTDE方法中性能最弱的。

\textbf{第三层：与MAPPO的差距——数值温和但结构重要。}本文方法相较MAPPO（当前最优CTDE-MARL基线）的改进幅度在数值上是温和的：冲突次数降低10.3\%（20.3$\to$18.2），延迟缩短12.3\%（0.65$\to$0.57 s），累计奖励提升44.2\%（892.4$\to$1286.7）。MAPPO作为强基线本身已具备相当的冲突规避能力——它通过集中式价值函数捕捉了部分全局协同信息。本文方法相较MAPPO的增量主要来自DyGAT的三个组件协同：（a）空间注意力比全连接网络更精准地聚焦高风险邻域飞行器；（b）时序注意力利用历史轨迹信息预判冲突趋势；（c）物理量驱动的动态边权重为注意力提供了有价值的先验引导。独立A2C由于完全缺乏智能体间显式协同机制，其在冲突次数（19.7次）上的表现虽与MAPPO（20.3次）接近甚至略优，但累计奖励（$-657.8$ vs.\ MAPPO的$+892.4$）和通行完成率（96.3\% vs.\ 98.2\%）大幅落后——这揭示了缺乏协同的核心代价：各飞行器为规避冲突采取过度保守的速度策略（频繁减速、保持过大间距），虽减少了冲突但严重牺牲了运行效率。独立A2C的低效率从反面验证了CTDE架构在兼顾安全与效率方面的必要性。

\subsubsection{安全指标专项对比——表层与深层指标的差异化图景}

表\ref{tab:safety_metrics}报告了四项安全攸关指标的对比结果。其中CFCR、MCAE与MSVR三项为本文引入的深层安全指标，揭示传统表层指标（TC、PCR）无法反映的安全质量信息。

\begin{table}[H]
    \centering
    \caption{安全指标专项对比（21架eVTOL，5次独立测试，均值$\pm$标准差）}
    \label{tab:safety_metrics}
    \begin{tabular}{|c|c|c|c|c|}
        \hline
        \textbf{方法} & \textbf{CFCR(\%)$\uparrow$} & \textbf{MCAE$\downarrow$} & \textbf{MSVR(\%)$\downarrow$} & \textbf{PCR(\%)$\uparrow$} \\
        \hline
        RR（传统规则） & $12.3 \pm 4.1$ & $2.03 \pm 0.25$ & $8.7 \pm 1.2$ & $89.2 \pm 3.5$ \\
        FCFS（传统规则） & $18.7 \pm 3.6$ & $1.85 \pm 0.22$ & $7.4 \pm 1.0$ & $91.5 \pm 2.8$ \\
        集中式DQN（单RL） & $37.5 \pm 3.2$ & $1.35 \pm 0.19$ & $5.1 \pm 0.8$ & $94.7 \pm 2.1$ \\
        独立A2C（独立ML） & $48.2 \pm 2.8$ & $0.94 \pm 0.13$ & $3.6 \pm 0.6$ & $96.3 \pm 1.7$ \\
        MADDPG（CTDE） & $42.6 \pm 3.0$ & $1.21 \pm 0.15$ & $4.2 \pm 0.7$ & $95.8 \pm 1.9$ \\
        COMA（CTDE） & $51.3 \pm 2.7$ & $1.05 \pm 0.14$ & $3.8 \pm 0.6$ & $97.1 \pm 1.5$ \\
        QMIX（CTDE） & $53.7 \pm 2.5$ & $1.10 \pm 0.13$ & $3.5 \pm 0.5$ & $97.5 \pm 1.3$ \\
        MAPPO（CTDE） & $57.8 \pm 2.2$ & $0.97 \pm 0.11$ & $3.0 \pm 0.4$ & $98.2 \pm 1.1$ \\
        \textbf{本文方法} & $\mathbf{68.4 \pm 1.8}$ & $\mathbf{0.87 \pm 0.10}$ & $\mathbf{2.4 \pm 0.3}$ & $\mathbf{99.1 \pm 0.8}$ \\
        \hline
    \end{tabular}
\end{table}

表\ref{tab:safety_metrics}揭示了三个层次的关键信息：

\textbf{第一层：PCR与CFCR的巨大鸿沟——表层指标掩盖的安全真相。}本文方法PCR $= 99.1\%$但CFCR仅$68.4\%$，两者相差30.7个百分点。这意味着约31\%的episode中，所有飞行器虽最终到达了目标区域（PCR计入），但途中至少经历了一次安全距离违反（CFCR不计入）。若仅报告PCR（如四月模型），读者将获得系统"几乎完美"的错误印象——这正是区分表层指标与深层指标、引入CFCR的核心方法论动机。MAPPO的PCR $= 98.2\%$与CFCR $= 57.8\%$之间也存在40.4个百分点的巨大差距，表明这一问题是CTDE-MARL方法的共性挑战，并非本文方法特有。

\textbf{第二层：MCAE与MSVR——安全质量的精细化度量。}MCAE $= 0.87$表明平均每架飞行器每episode经历不到1次冲突，较MAPPO的0.97次降低了10.3\%，与TC的降幅一致。MSVR $= 2.4\%$意味着仅2.4\%的仿真时间步存在安全距离违反——即97.6\%的时间步中系统处于安全状态，较MAPPO的3.0\%降低了20.0\%（相对降幅）。MSVR的综合降幅大于MCAE，表明本文方法不仅减少了冲突的频次（数量维度），更缩短了每次冲突的持续时间（时间维度）——飞行器在发生安全距离违反后能够更快地恢复安全间隔。

\textbf{第三层：规则方法与学习方法的CFCR鸿沟。}RR的CFCR仅为12.3\%，PCR为89.2\%——意味着仅约1/8的episode实现了全程无冲突通行，但近90\%的飞行器最终仍到达了目标。这一极端不对称揭示了规则方法的工作模式：飞行器在密集冲突中"跌跌撞撞"通过，而非协调有序通过。学习方法的CFCR系统性地高于规则方法——这是学习方法能够主动协调飞行器航迹的直接证据。

\subsubsection{统计显著性检验}

对本文方法与MAPPO在四项关键指标上的性能差异进行Welch双样本$t$检验（双侧，不等方差，$n_1 = n_2 = 5$），结果如表\ref{tab:statistical_tests}所示。选择Welch $t$检验而非Student $t$检验的原因为：两种方法的性能方差不等（Levene检验$p < 0.05$），Welch检验对方差齐性不敏感，更为保守和可靠。

\begin{table}[H]
    \centering
    \caption{本文方法 vs.\ MAPPO的统计显著性检验（$n=5$，双侧Welch $t$检验，$\alpha=0.05$）}
    \label{tab:statistical_tests}
    \begin{tabular}{|c|c|c|c|c|c|}
        \hline
        \textbf{指标} & \textbf{均值差} & \textbf{$t$统计量} & \textbf{$p$值} & \textbf{Cohen's $d$} & \textbf{显著性} \\
        \hline
        TC（总冲突次数） & $2.1 \downarrow$ & $t=1.47$ & $p=0.180$ & $0.93$ & 不显著（n.s.） \\
        \hline
        CFCR（无冲突完成率） & $10.6\% \uparrow$ & $t=8.34$ & $p<0.001$ & $5.27$ & \textbf{高度显著} \\
        \hline
        MCAE（每架次冲突） & $0.10 \downarrow$ & $t=1.50$ & $p=0.172$ & $0.95$ & 不显著（n.s.） \\
        \hline
        MSVR（最小间隔违反率） & $0.6\% \downarrow$ & $t=2.68$ & $p=0.029$ & $1.69$ & \textbf{显著} \\
        \hline
    \end{tabular}
\end{table}

表\ref{tab:statistical_tests}揭示了一个细微但关键的统计图景，需要谨慎解读以避免两类错误：

\textbf{（1）TC的不显著性不应等同于"两方法无差异"。}TC差异$p=0.180$在$\alpha=0.05$水平上未达统计显著，但这$\neq$"本文方法与MAPPO在TC上无差异"。事后统计功效分析表明：Cohen's $d \approx 0.93$属于大效应量（$>0.8$），但在$n=5$的小样本下，双侧Welch $t$检验的统计功效（power）仅约38\%——即有62\%的概率，即使真实效应存在，我们也无法在$\alpha=0.05$水平上检测到它。达到80\%统计功效（$\alpha=0.05$，双侧）以检测$d=0.93$的效应量，约需每组$n \geq 20$的样本量。因此，TC的"不显著"最恰当的解读是：\textbf{现有样本量下的统计功效不足，无法在$\alpha=0.05$水平上确认TC差异}——而非\textbf{确认TC无差异}。这是两类统计推断错误的典型案例。

\textbf{（2）CFCR的极显著差异——效应量大到小样本也足以检测。}CFCR差异的$p<0.001$（$t=8.34$，Cohen's $d \approx 5.27$，超大效应量）即使在小样本（$n=5$）下也达到$>99\%$的统计功效——效应量如此之大，5次独立测试足以提供压倒性的证据。这一结论高度稳健，对$n$的增加或$\alpha$的调整不敏感。

\textbf{（3）指标间的差异化统计图景揭示了本文方法的真实增量性质。}TC不显著但CFCR极显著的模式表明：本文方法相较MAPPO的优势不在于大幅降低冲突的绝对频次（所有CTDE方法在此维度已表现良好——MAPPO本身已有$20.3$次/episode的冲突水平），而在于将更多的episode从"有冲突通过"转化为"全程无冲突通过"——即提升了系统的\textbf{安全一致性}（safety consistency）。DyGAT的细粒度时空建模使飞行器能够更精准地协调航迹，减少冲突的"残余"——那些MAPPO未能规避但在整体统计中被平均效应掩盖的冲突。

\textbf{（4）MSVR的显著差异——时间维度的安全改善。}MSVR差异$p=0.029$（$d \approx 1.69$，大效应量）在$\alpha=0.05$水平上显著，表明本文方法在缩短冲突持续时间维度上同样具有统计可信的改善。这为DyGAT的时序注意力机制提供了间接但有力的实证支持——时序注意力使飞行器能够在安全距离违反发生后更快地恢复安全间隔。

\subsection{消融实验}

表\ref{tab:ablation}报告了各核心模块的独立性能贡献。消融实验遵循"一次移除一个模块，其余保持不变"的单变量控制原则，所有消融变体均使用相同的随机种子集进行5次独立测试。

\begin{table}[H]
    \centering
    \caption{消融实验结果（21架eVTOL，5次独立测试，均值$\pm$标准差）}
    \label{tab:ablation}
    \resizebox{0.92\textwidth}{!}{%
    \begin{tabular}{|c|c|c|c|c|}
        \hline
        \textbf{模型配置} & \textbf{累计奖励$\uparrow$} & \textbf{冲突次数$\downarrow$} & \textbf{平均延迟(s)$\downarrow$} & \textbf{PCR(\%)$\uparrow$} \\
        \hline
        完整模型（本文） & $\mathbf{1286.7 \pm 18.2}$ & $\mathbf{18.2 \pm 2.1}$ & $\mathbf{0.57 \pm 0.03}$ & $\mathbf{99.1 \pm 0.8}$ \\
        \hline
        $-$ DyGAT（改用标准GAT） & $215.3 \pm 32.6$ & $49.8 \pm 4.5$ & $0.79 \pm 0.05$ & $94.3 \pm 1.9$ \\
        \hline
        $-$ 课程学习（直接$N=21$训练） & $587.2 \pm 27.4$ & $28.4 \pm 3.1$ & $0.68 \pm 0.04$ & $97.2 \pm 1.4$ \\
        \hline
        $-$ 冲突轨迹加权（$\lambda \equiv 1.0$） & $763.5 \pm 24.1$ & $24.6 \pm 2.8$ & $0.65 \pm 0.04$ & $97.8 \pm 1.2$ \\
        \hline
        $-$ 课程学习 $-$ 冲突轨迹加权 & $128.6 \pm 38.5$ & $35.7 \pm 3.9$ & $0.82 \pm 0.06$ & $95.6 \pm 1.7$ \\
        \hline
    \end{tabular}%
    }
\end{table}

\textbf{DyGAT是性能贡献最大的单一模块。}移除DyGAT（替换为标准GAT——仅空间GAT注意力，无时序注意力、无动态边权重）后，冲突次数从18.2激增至49.8，相对增幅达173.6\%，累计奖励从1286.7骤降至215.3。这一幅度远超其他任何模块的单独移除效果（课程学习56.0\%，轨迹加权35.2\%），确定性地表明DyGAT的时空耦合建模是本文方法性能的\textbf{绝对基石}。更深层的解读是：GAT的空间注意力本身不足以应对高动态场景——正是时序注意力（预判运动趋势）与动态边权重（区分风险等级）的协同，才使DyGAT超越了标准GAT。

\textbf{课程学习与冲突轨迹加权形成互补协同。}单独移除课程学习（直接以$N=21$全程训练）导致冲突上升56.0\%，而单独移除轨迹加权（$\lambda \equiv 1.0$）导致冲突上升35.2\%。两者同时移除时冲突上升96.2\%（18.2$\to$35.7），该增幅与单独效应之和（56.0\%$+$35.2\%$=91.2\%$）在统计误差范围内基本一致（5个百分点的差异在$n=5$下不显著），表明两种机制的作用路径大部分相互独立——课程学习主要解决"场景难度适配"问题（从简到繁的渐进训练），轨迹加权主要解决"样本效率"问题（在每一难度阶段放大冲突样本的损失贡献）。

\textbf{当前消融实验的粒度局限：}需要坦诚指出两项不足，并将其列为投稿前的优先补充实验。（1）DyGAT的移除是整体替换为标准GAT，尚未分解为空间注意力、时序注意力与动态边权重的独立消融——173.6\%的退化是三者的联合效应，各子模块的独立贡献量未知。投稿前计划补充的三组独立消融为：仅移除时序注意力（保留空间GAT$+$动态边权重）、仅移除动态边权重（保留空间GAT$+$时序注意力）、仅使用空间GAT（baseline for comparison）。（2）$\lambda=2.0$仅与$\lambda=1.0$对比，其对$\lambda \in \{1.5, 2.5, 3.0\}$的敏感性未经充分验证——$\lambda=2.0$的最优性无法确认。

\subsection{密度鲁棒性与敏感性分析}

\subsubsection{飞行器密度鲁棒性}

图\ref{fig:density_robustness}展示了$N \in \{5, 10, 15, 21\}$四种密度场景下的冲突次数对比。本文方法的冲突次数增长斜率（$\approx 0.75$次/架）显著低于MAPPO（$\approx 1.15$次/架）与RR（$\approx 2.4$次/架）。在低密度场景（$N=5$）中，所有CTDE-MARL方法的冲突次数均处于较低水平（本文2.1次，MAPPO 2.4次，QMIX 2.8次）——方法间差异小，因为低密度下冲突本身稀少，精细建模的边际收益有限。密度增至$N=15$时，本文方法相较MAPPO的冲突优势扩大至22\%（14.3 vs.\ 18.4次），密度增至$N=21$时为10.3\%。非线性趋势（15架$\to$21架时优势略有收窄）的可能原因为：在极高密度下，物理可行解空间被压缩到极限——即使完美的时空建模也无法完全规避冲突，此时环境密度成为冲突次数的硬约束。

DyGAT在密度升高时优势愈发显著的核心机制在于：飞行器越多，邻域内飞行器对越多（$N=21$时最大邻域内可达20架），精准区分不同飞行器对的冲突风险等级越关键——将有限的注意力资源集中于高风险交互对，而非均匀分配。

\begin{figure}[H]
    \centering
    \includegraphics[width=0.65\textwidth]{QQ图片20260524171218.png}
    \caption{不同飞行器密度（$N=5, 10, 15, 21$）下的冲突次数对比。本文方法、MAPPO与RR三类方法分别展示。误差条表示$\pm1$标准差（5次独立测试）。}
    \label{fig:density_robustness}
\end{figure}

\subsubsection{超参数敏感性分析}

图\ref{fig:sensitivity_analysis}在$N=21$场景下测试了三类关键参数的影响：

\textbf{基础安全距离 $d_{\text{base}}$（40--70 m）：}$d_{\text{base}}$从40 m增至70 m时，冲突次数从22.5降至14.7（$\downarrow 34.7\%$），但延迟从0.52 s升至0.63 s（$\uparrow 21.2\%$），体现了航空交通管理中经典的安全-效率基本权衡：更大的安全距离降低冲突概率（安全收益），但限制了飞行器的速度优化空间（效率代价）。操作点$d_{\text{base}}=50$ m在当前奖励设计下位于该权衡曲线的"膝点"附近。

\textbf{风扰强度（0--0.3 m/s）：}风扰从0增至0.3 m/s时，冲突次数上升27.6\%（从约16.0增至约20.4）。性能退化幅度温和，DyGAT的时序注意力机制通过学习历史轨迹的时序趋势，部分补偿了风扰带来的位置预测不确定性。但需要着重指出：当前均匀随机风扰模型（$\theta_w \sim \mathcal{U}(0, 2\pi)$，固定强度，时间独立）是对真实低空湍流的显著简化。Dryden风湍流模型（MIL-F-8785C）\cite{ref20}以连续时间随机过程描述风场速度分量，具有特定的功率谱密度（von K\'{a}rm\'{a}n谱）与空间相干结构——这些特性在均匀随机风扰中完全缺失。因此，0.3 m/s强度下仅27.6\%的性能退化可能部分归因于风扰模型本身过于"友好"（独立同分布的随机扰动易于被时序平滑抵消），真实Dryden风场下的性能退化幅度可能更大——这一假设需要在未来工作中使用标准风谱模型进行验证。

\textbf{传感器噪声 $\sigma$（0--1.5 m）：}$\sigma$从0增至1.5 m时，冲突次数仅上升14.6\%——所有参数中影响最小。这一低敏感性的可能原因有二：（a）LSTM编码器对连续10步（5 s）历史状态序列的隐式平滑作用有效过滤了高频观测噪声；（b）$\sigma=1.5$ m相对$d_{\text{safety}} \approx 85$ m（$N=21$时典型值）的比例仅为1.8\%，仍处于安全距离的感知容差范围内。

\begin{figure}[H]
    \centering
    \includegraphics[width=0.8\textwidth]{QQ图片20260524162300.png}
    \caption{超参数敏感性分析——分别改变基础安全距离$d_{\text{base}}$（左）、风扰强度（中）与传感器噪声$\sigma$（右）时冲突次数与平均延迟的变化。每个数据点为5次独立测试的均值，误差条为$\pm1$标准差。红色虚线标记默认参数值。}
    \label{fig:sensitivity_analysis}
\end{figure}

\subsection{计算成本分析}

图\ref{fig:computation_cost}展示了不同$N$下各方法在NVIDIA RTX 4060 GPU上的单步推理平均耗时（含状态预处理、网络前向传播与动作后处理，不含环境步进时间）。本文方法在$N=21$时平均推理时间为$4.5$ ms，远低于仿真步长$500$ ms——实时性裕度约为111$\times$，满足实时部署的时序约束。

计算时间从$N=5$时的$1.2$ ms增长至$N=21$时的$4.5$ ms，增长斜率约为$0.21$ ms/架，在所有CTDE-MARL方法中最低。这一优异的计算可扩展性得益于DyGAT的两项设计特征：（1）动态边权重$w_{ij}$基于物理量（距离、速度差、航向差）直接计算，无可学习参数，不涉及梯度回传或矩阵乘法——每对飞行器的$w_{ij}$计算仅需3次指数运算和2次乘法；（2）邻域限制（500 m）将每个飞行器的注意力计算限制在局部范围内，而非全图——最坏情况下（$N=21$且所有飞行器集中在500 m邻域内），每个飞行器仅需计算20个注意力权重。

COMA在$N=21$时耗时超过$15$ ms（反事实基线需为每个智能体分别计算边缘化期望，涉及$N$次独立的Critic前向传播与策略采样），QMIX从$2.1$ ms（$N=5$）增至约$8$ ms（$N=21$）（超网络输入维度随$N$线性增长）。

\begin{figure}[H]
    \centering
    \includegraphics[width=0.8\textwidth]{QQ图片20260524163903.jpg}
    \caption{不同飞行器数量下的平均单步推理计算时间对比（NVIDIA RTX 4060 GPU）。本文方法、MAPPO、QMIX与COMA四类CTDE-MARL方法展示。嵌入小图放大展示$N=5\sim10$区间的差异。}
    \label{fig:computation_cost}
\end{figure}

%==========================================================================
% 5. 讨论
%==========================================================================
\section{讨论}

\subsection{核心贡献的机制解读与功能归因}

本文方法的性能优势可归因于三个设计要素的协同，每个要素的功能定位与贡献量级不同。

\textbf{DyGAT的精准时空耦合刻画是性能基石。}消融实验中DyGAT移除导致冲突暴增173.6\%，其贡献量级远超其他模块的总和。这一定量发现内含一个重要的定性启示：在低空交通管控中，\textbf{图结构的质量是冲突风险的主导因素}——它直接决定模型对邻域风险的感知精度。一个值得注意的不对称性：移除DyGAT后冲突暴增173.6\%，而$d_{\text{base}}$从40 m变化至70 m时冲突仅变化34.7\%。这揭示了二者的功能定位差异：DyGAT决定"如何感知和规避冲突"（方法层面的创新），安全距离定义"什么是冲突"（环境层面的运营参数）。两者不应被等同视之——安全距离的工程改进应定位为环境设计贡献，不应与DyGAT的方法贡献混淆。

\textbf{课程学习与冲突轨迹加权形成训练机制协同。}课程学习提供从简到繁的渐进训练路径（$5 \to 21$架，$d_{\text{base}}$从80 m递减至50 m），轨迹加权（$\lambda=2.0$）确保每一难度阶段中稀疏的冲突样本被充分学习。两者单独移除的效应（56.0\%和35.2\%）之和约等于同时移除的效应（96.2\%），表明两种机制的作用路径高度独立但最终目标一致——提升训练稳定性与样本效率。

\textbf{GAT-Mixer的架构适配性。}完全图GAT聚合天然适配$N$动态变化的课程学习场景（无需重新设计网络结构），且注意力权重的无约束特性突破了QMIX的单调性限制。但GAT-Mixer的性能贡献是间接的——它通过改善训练信号的质量（更准确的全局价值估计）间接提升策略质量，而非直接增强策略的感知或决策能力。

\subsection{CFCR的深层工程含义与安全-效率权衡}

CFCR $= 68.4\%$是全文最值得深入讨论的数字，需要从多个视角进行审慎分析。

\textbf{乐观视角：}本文方法在约三分之二的测试episode中实现了全程零冲突安全通行——即这些episode中所有$N=21$架飞行器在100 s仿真时间内从未违反安全距离约束。CFCR较MAPPO的57.8\%提升10.6个百分点，且该差异在统计上极显著（$p<0.001$，Cohen's $d \approx 5.27$），具有高度的结论稳健性。从绝对数值看，68.4\%显著高于其他任何基线方法（第二高的MAPPO为57.8\%）。

\textbf{悲观视角：}仍有约31\%的episode存在至少一次安全距离违反。在航空安全语境下，一个关键问题是：31\%的"有冲突通过"episode中冲突的严重程度分布如何？当前实验尚未按冲突严重程度进行分层统计——这是本文实验设计最重要的遗留缺陷。

为解决这一问题，建议将安全距离违反按$d_{\min}/d_{\text{safety}}$比值分为三个严重程度等级：

\begin{itemize}
    \item \textbf{轻微违规（0.8 $\leq d_{\min}/d_{\text{safety}} < 1.0$）：}瞬时、小幅度的安全距离偏离，通常由传感器噪声或风扰导致的短时超调引起。在工程实践中，此类违规在低空交通中可能在一定容差范围内被容忍。
    \item \textbf{中等违规（0.5 $\leq d_{\min}/d_{\text{safety}} < 0.8$）：}显著的安全距离违反，表明飞行器间存在真实的碰撞风险，但通过紧急机动仍可矫正。
    \item \textbf{严重违规/NMAC（$d_{\min}/d_{\text{safety}} < 0.5$）：}接近空中碰撞（Near Mid-Air Collision），在任何运营标准下均不可接受。国际民航组织（ICAO）将NMAC定义为飞行器间距离小于500 ft（约152 m）；在本文场景中对应$d_{\min}/d_{\text{safety}} < 0.5$（$d_{\text{safety}} \approx 85$ m时，$<42.5$ m）。
\end{itemize}

CFCR $= 68.4\%$的可接受性关键取决于31\%的"有冲突通过"episode在这三个层级上的分布：若90\%以上为轻微违规（可容忍的瞬时偏离），68.4\%的CFCR可以被辩护——系统在绝大多数情况下是安全的，偶发的轻微超调在工程上可接受；若包含显著比例（$>5\%$）的严重违规/NMAC级别事件，则系统在当前形态下\textbf{不安全}——即使CFCR达到99\%，剩余的1\%若包含NMAC，系统仍然不可接受。

参考民航终端区的碰撞风险目标（每飞行小时$10^{-7}$的目标安全水平），CFCR在部署前至少应达到95\%以上（且剩余5\%中NMAC级别事件应低于$10^{-5}$每飞行小时）。从68.4\%到95\%+的提升路径包括：
\begin{enumerate}[label=(\arabic*)]
    \item 引入ORCA/CBF作为底层安全网，构成"MARL高层优化$+$形式化方法底层安全保障"的双层架构——MARL负责效率优化与协同协调，CBF负责硬安全约束的实时执行；
    \item 在奖励函数中增大严重违规（$d_{\min}/d_{\text{safety}} < 0.5$）的惩罚梯度，使策略对NMAC级别事件具有更强的敏感性；
    \item 延长训练轮次（当前1000轮可能不足以收敛至最优安全策略）并增加初始位置与目标速度的随机化程度，提升策略对边缘情况的覆盖；
    \item 采用标准风谱模型（Dryden/Kaimal）进行更真实的鲁棒性验证，确保CFCR在真实风场下不显著退化。
\end{enumerate}

\textbf{投稿前优先补充：}对已有5次独立测试（每次测试包含多个episode）的运行日志中所有安全距离违反事件进行$d_{\min}/d_{\text{safety}}$直方图统计，并按三级严重程度分层报告分布比例——这是支撑CFCR相关讨论的实证基础。

\subsection{与MS-GRL的详细比较}

Li等提出的MS-GRL\cite{ref24}与本文目标最为接近——同样面向密集低空网络的冲突解决，同样采用图注意力机制。表\ref{tab:ms_grl_comparison}从四个维度给出系统性对比。

\begin{table}[H]
    \centering
    \caption{本文方法与MS-GRL的系统性对比}
    \label{tab:ms_grl_comparison}
    \begin{tabular}{|c|c|c|}
        \hline
        \textbf{对比维度} & \textbf{MS-GRL \cite{ref24}} & \textbf{本文方法} \\
        \hline
        算法选择 & Double DQN（离散动作：速度档位选择） & PPO（连续动作：加速度控制） \\
        \hline
        图结构 & 多尺度GAT，固定距离阈值 & DyGAT（空间$+$时序注意力$+$物理量动态边权重） \\
        \hline
        场景规模 & 通用UAV，100架 & eVTOL交叉路口，21架 \\
        \hline
        动作空间 & 离散（$k$个速度档位） & 连续（$a \in \mathbb{R}$，加速度m/s$^2$） \\
        \hline
        安全约束 & 固定安全距离 & 与速度正相关的动态安全距离 \\
        \hline
        训练策略 & 标准经验回放 & 课程学习$+$on-policy轨迹加权 \\
        \hline
    \end{tabular}
\end{table}

两者在技术路线上的本质差异决定了互补而非替代的关系：

\textbf{（1）算法选择（Double DQN vs.\ PPO）。}MS-GRL采用Double DQN处理离散动作空间——将速度调控量化为有限个档位（如加速/减速/保持）。离散动作的优势在于训练稳定（避免连续动作空间的高方差梯度估计），代价是控制粒度粗糙——飞行器无法根据冲突风险的连续变化做出精细的速度调节。本文采用PPO处理连续加速度控制，允许飞行器在物理约束范围内任意调节加速度——这在eVTOL的精细化速度调控需求中是更自然的选择，但训练难度更大（连续动作空间的探索效率更低）。

\textbf{（2）图结构（固定阈值 vs.\ 物理量驱动动态权重）。}MS-GRL的多尺度GAT图结构基于固定距离阈值——例如，分别在100 m、300 m和500 m三个尺度上构建子图，聚合各尺度的特征。固定阈值无法区分"距离相同但相对运动状态不同"的飞行器对的冲突风险差异。本文的DyGAT通过距离、速度差和航向差三个物理量计算连续动态边权重$w_{ij} \in (0, 1]$——两对距离同为200 m的飞行器，若一对同向等速飞行（$w_{ij} \approx 1.0$），另一对相向高速靠近（$w_{ij} \approx e^{-0.2} \cdot e^{-4} \cdot e^{-0} \approx 0.015$），后者的注意力权重仅为前者的约1/67，风险区分能力远超固定阈值方案。

\textbf{（3）场景规模与针对性（通用UAV vs.\ 特定运营场景）。}MS-GRL面向通用UAV编队场景（100架规模），强调大规模可扩展性。本文聚焦eVTOL交叉路口的特定运营场景（21架规模），在安全约束（动态安全距离）、物理建模（eVTOL运动学约束）与运营逻辑（交叉路口通行规则）上更具针对性。

\subsection{方法局限性}

本文方法存在以下主要局限性，按严重程度排序：

\textbf{（1）场景维度限制（严重程度：高）。}当前仿真场景为二维单交叉路口的简化模型。真实低空交通的本质特征包括：三维高度层（含垂直起降与巡航高度切换）、多路口网络拓扑（交叉口之间的交通流耦合）、起降场容量约束（垂直起降阶段的空中等待）与通航走廊的空间约束。二维简化忽略了垂直维度的冲突规避——在真实三维空间中，飞行器可通过高度层分离规避冲突，而无需仅依赖水平面的速度/航向调整。这是当前模型最根本的场景限制，限制了结论向真实三维广域低空交通的外推有效性。

\textbf{（2）通信假设的理想化（严重程度：高）。}CTDE训练假设全局状态$S$可在集中式训练服务器上无延迟、无丢失地获取。在真实UTM部署环境中，飞行器与地面站之间的通信链路面临延迟（典型值50--200 ms）、数据包丢失（城市环境多径效应）与有限带宽（多飞行器共享频谱资源）的约束。通信受限将同时影响两个层面：训练阶段——全局状态的不完整/延迟获取导致GAT-Mixer的价值估计偏差；执行阶段——虽然各飞行器独立决策不依赖通信，但缺乏全局信息的策略可能在通信中断时做出次优决策。

\textbf{（3）飞行器同质化假设（严重程度：中）。}所有$N$架飞行器使用相同的运动学参数（加速度、减速度、速度范围、转弯角约束）。实际AAM生态中，不同制造商（Joby、Volocopter、Lilium、Archer、亿航等）与载荷等级（2座、4座、6座）的eVTOL具有差异化的飞行包线——最大巡航速度可能从80 km/h（轻型多旋翼）到320 km/h（倾转旋翼构型）不等，制动性能差异可达数倍。同质化假设简化了问题，但限制了模型在异构飞行器群体中的直接适用性。

\textbf{（4）风扰与噪声建模的简化（严重程度：中）。}风扰采用均匀随机分布（$\theta_w \sim \mathcal{U}(0, 2\pi)$，固定强度0.1 m/s，时间独立），噪声采用高斯白噪声（$\mathcal{N}(0, 0.5^2)$，时间空间独立）。真实低空风场具有显著的空间相关性与时间连续性——Dryden模型（MIL-F-8785C）\cite{ref20}基于von K\'{a}rm\'{a}n湍流谱描述连续时间随机风场，其湍流积分尺度（$L_u, L_v, L_w$）可达数百米，意味着空间上邻近的飞行器经历的风扰是统计相关的（而非独立的）。第4.5.2节的敏感性分析结果——风扰从0增至0.3时的27.6\%性能退化——可能低估了真实Dryden风场下因空间相干结构导致的性能退化。

\textbf{（5）基线比较范围不足（严重程度：中）。}当前实验未纳入MPC\cite{ref27}、ORCA/RVO\cite{ref22}、CBF\cite{ref23}及基于MILP的集中式时空资源分配方法作为基线。这些方法在形式化安全保证（CBF确保状态始终在安全集内）、最优性证明（MPC在模型精确时提供有限时域最优解）或计算效率（ORCA在100架以上仍可实时运行）方面具有独特优势，与本文方法形成互补。当前仅战胜RR、FCFS、DQN和几种CTDE-MARL方法，不足以构成完整的贡献声明——与这些工程安全关键方法的比较是投稿前或修改阶段的优先补充实验。

\textbf{（6）实验分析的粒度局限（严重程度：中低）。}三项具体不足：（a）冲突严重程度分布（$d_{\min}/d_{\text{safety}}$直方图及三级分层统计）未报告；（b）DyGAT的组件级消融（空间注意力、时序注意力、动态边权重的独立贡献分离）未完成；（c）$\lambda$超参数仅测试了$\{1.0, 2.0\}$两个取值，最优区间未经验证；（d）时序融合系数$\alpha_{\text{temp}}=0.3$的敏感性分析（测试值$\{0.1, 0.2, 0.3, 0.5, 0.7\}$）的粒度较粗，最优值及其置信区间未报告。

\subsection{实际应用潜力与未来工作}

在充分认识第5.4节所列局限性的前提下，本文模型的合理定位是：作为低空UTM冲突管理模块的\textbf{候选算法}，为高密度eVTOL自主协同管控提供技术储备——而非声称可直接部署至生产系统。分布式执行架构使计算负载分散至各飞行器机载设备（$N=21$时4.5 ms/步的推理耗时远低于500 ms仿真步长，实时性裕度$\approx 111\times$），避免了集中式决策的单点故障风险，是分布式架构在实际部署中的重要优势。

\textbf{投稿前优先补充实验（短期，1--2个月）：}
\begin{enumerate}[label=(\arabic*)]
    \item 对已有测试运行日志进行冲突严重程度分布（$d_{\min}/d_{\text{safety}}$直方图）统计，按三级分层报告；
    \item 完成DyGAT组件级消融：独立移除空间注意力、时序注意力与动态边权重，量化各子模块的独立贡献；
    \item $\lambda \in \{1.0, 1.5, 2.0, 2.5, 3.0\}$的参数敏感性分析与最优区间确定；
    \item $\alpha_{\text{temp}} \in \{0.1, 0.15, 0.2, 0.25, 0.3, 0.4, 0.5\}$的精细网格搜索，验证0.3的最优性。
\end{enumerate}

\textbf{长期研究方向（中期至长期，6--24个月）：}
\begin{enumerate}[label=(\arabic*)]
    \item \textbf{三维广域空域扩展：}将仿真环境从二维单交叉路口扩展至三维多路口网络，引入垂直起降动力学（推力矢量、过渡阶段）与起降场容量约束；
    \item \textbf{工程安全基线集成：}与MPC、ORCA/RVO、CBF等方法进行系统基准对比，探索"MARL优化$+$形式化安全保障"的混合架构；
    \item \textbf{标准风场验证：}采用Dryden（MIL-F-8785C）与Kaimal（IEC 61400-1）标准风谱模型替代当前均匀随机风扰，进行更真实的鲁棒性评估；
    \item \textbf{通信受限鲁棒训练：}引入通信延迟、丢包与有限带宽约束，研究通信降级条件下CTDE训练与分布式执行的性能退化及缓解策略；
    \item \textbf{异构飞行器群体：}引入差异化的飞行器动力学参数（速度范围、加速度、转弯半径），研究异构群体中的策略泛化与自适应；
    \item \textbf{硬件在环仿真与实飞验证：}在PX4/ArduPilot等飞控硬件在环环境中验证策略的实时性，最终进行小规模（3--5架）实飞验证。
\end{enumerate}

代码将在论文接收后开源，以促进低空MARL管控领域的可复现研究。

%==========================================================================
% 6. 结论
%==========================================================================
\section{结论}

本文针对城市低空十字交叉路口的eVTOL协同管控与冲突避免问题，提出了一种基于耦合多智能体强化学习的管控模型。模型核心设计——DyGAT动态图注意力网络、课程学习与冲突轨迹加权训练框架——分别在时空耦合建模与稀疏奖励处理两个关键维度上做出了针对性改进。

在21架eVTOL交叉路口仿真场景中，本文方法取得了以下主要实验结果：
\begin{itemize}
    \item 相较规则基线（RR），冲突次数降低57.2\%，平均延迟缩短63.9\%——验证了学习方法相较固定规则的结构性优势；
    \item 相较当前最优CTDE-MARL方法MAPPO，冲突次数降低10.3\%（$p=0.180$，未达统计显著），但无冲突完成率从57.8\%提升至68.4\%（$p<0.001$），最小间隔违反率从3.0\%降至2.4\%（$p=0.029$）——验证了方法在提升安全一致性方面的真实增量；
    \item 消融实验表明DyGAT是性能贡献最大的单一模块（移除后冲突增加173.6\%），课程学习与冲突轨迹加权形成互补协同。
\end{itemize}

本文的核心方法论贡献在于：（1）引入CFCR等深层安全指标，揭示了传统表层指标（PCR）掩盖的安全质量差异——PCR $=99.1\%$与CFCR $=68.4\%$之间30.7个百分点的鸿沟是全文最重要的实证发现；（2）通过Welch $t$检验与事后统计功效分析，诚实呈现了TC不显著（$p=0.180$，功效仅38\%）与CFCR高度显著（$p<0.001$，功效$>99\%$）的差异化统计图景——这种区分统计功效不足与效应量为零的分析范式，对于安全攸关的AI系统评估具有方法论示范价值。

本文的局限性（二维简化场景、均匀风扰模型、有限基线覆盖、$n=5$小样本量、冲突严重程度分布未统计）在第5.4节中已有坦诚讨论。这些局限意味着当前结果应被解读为概念验证（proof-of-concept），而非生产就绪（production-ready）系统的性能声明。投稿前的优先补充实验（DyGAT组件级消融、$\lambda$参数敏感性、冲突严重程度分层统计）是支撑本文核心结论的必要实证基础。

在充分认识上述局限的前提下，本文工作为城市低空交通的分布式自主管控提供了一条有潜力的技术路线——其核心设计理念（物理量驱动的图结构动态权重、on-policy框架下的稀疏样本轨迹加权、安全一致性的深层评价）对更广泛的低空MARL管控研究具有方法论参考价值。

%==========================================================================
% 参考文献
%==========================================================================
\begin{thebibliography}{99}

% [1] FAA — AAM预测报告 (首个引用)
\bibitem{ref26} Federal Aviation Administration. Urban air mobility (UAM): concept of operations v2.0[R]. Washington, DC: FAA, 2023.

% [2] Menon等 — 交通流欧拉模型ATFM (JGCD 2004)
\bibitem{ref31} MENON P K, SWERIDUK G D, BILIMORIA K D. New approach to modeling, analysis, and control of air traffic flow[J]. Journal of Guidance, Control, and Dynamics, 2004, 27(5): 737-744.

% [3] MPC — 模型预测控制航空应用 (ACC 2002)
\bibitem{ref27} RICHARDS A, HOW J P. Aircraft trajectory planning with collision avoidance using mixed integer linear programming[C]//Proceedings of the 2002 American Control Conference (ACC). Piscataway: IEEE, 2002: 1936-1941.

% [4] ORCA — 速度障碍法 (ISRR 2011)
\bibitem{ref22} VAN DEN BERG J, GUY S J, LIN M, et al. Reciprocal n-body collision avoidance[C]//Robotics Research: The 14th International Symposium ISRR (Springer Tracts in Advanced Robotics, Vol. 70). Berlin: Springer, 2011: 3-19.

% [5] CBF — 控制障碍函数综述 (ECC 2019)
\bibitem{ref23} AMES A D, COOGAN S, EGERSTEDT M, et al. Control barrier functions: theory and applications[C]//2019 18th European Control Conference (ECC). Piscataway: IEEE, 2019: 3420-3431.

% [6] Agogino & Tumer — 耦合智能体调控框架 (AAMAS 2008)
\bibitem{ref1} AGOGINO A, TUMER K. Regulating air traffic flow with coupled agents[C]//Proceedings of the 7th International Conference on Autonomous Agents and Multiagent Systems (AAMAS 2008). Richland: IFAAMAS, 2008: 535-542.

% [7] Agogino & Tumer — 多智能体空中交通流量管理 (JAAMAS 2012)
\bibitem{ref2} AGOGINO A K, TUMER K. A multiagent approach to managing air traffic flow[J]. Autonomous Agents and Multi-Agent Systems, 2012, 24(1): 1-25.

% [8] Brittain & Wei — 自主间隔保持 (ITSC 2019)
\bibitem{ref3} BRITTAIN M, WEI P. Autonomous separation assurance in a high-density en route sector: a deep multi-agent reinforcement learning approach[C]//2019 IEEE Intelligent Transportation Systems Conference (ITSC). Piscataway: IEEE, 2019: 3256-3262.

% [9] Ghosh等 — 深度集成MARL (ICAPS 2021)
\bibitem{ref9} GHOSH S, LAGUNA S, LIM S H, et al. A deep ensemble method for multi-agent reinforcement learning: a case study on air traffic control[C]//Proceedings of the 31st International Conference on Automated Planning and Scheduling (ICAPS 2021). Palo Alto: AAAI Press, 2021: 468-476.

% [10] Xie等 — RL流量管理 (DASC 2021)
\bibitem{ref6} XIE Y B, GARDI A, SABATINI R. Reinforcement learning-based flow management techniques for urban air mobility and dense low-altitude air traffic operations[C]//2021 IEEE/AIAA 40th Digital Avionics Systems Conference (DASC). Piscataway: IEEE, 2021: 1-10.

% [11] MS-GRL — 多尺度图强化学习无人机冲突解决 (IEEE IoT-J 2025)
\bibitem{ref24} LI X, WANG Y, ZHANG H, et al. Multiscale-graph-enhanced reinforcement learning for conflict resolution in dense UAV networks[J]. IEEE Internet of Things Journal, 2025, 12(22): 72134-72148.

% [12] Spatharis等 — 层级MARL空管方案 (NCA 2021)
\bibitem{ref7} SPATHARIS C, BASTAS A, KRAVARIS T, et al. Hierarchical multiagent reinforcement learning schemes for air traffic management[J]. Neural Computing and Applications, 2021, 33(14): 8649-8671.

% [13] Aslani等 — Actor-Critic交通信号 (TR_C 2017)
\bibitem{ref10} ASLANI M, MESGARI M S, WIERING M. Adaptive traffic signal control with actor-critic methods in a real-world traffic network with different traffic disruption events[J]. Transportation Research Part C: Emerging Technologies, 2017, 85: 732-752.

% [14] 翟子洋等 — RL/DRL交通信号综述
\bibitem{ref11} 翟子洋, 郝茹茹, 董世浩. 大规模智慧交通信号控制中的强化学习和深度强化学习方法综述[J]. 计算机应用研究, 2024, 41(6): 1618-1627.

% [15] Dryden风湍流模型标准 (MIL-F-8785C)
\bibitem{ref20} U.S. Department of Defense. Flying qualities of piloted airplanes: MIL-F-8785C[S]. Washington, DC: U.S. Department of Defense, 1980.

% [16] Tumer & Agogino — 分布式智能体流量管理 (AAMAS 2007)
\bibitem{ref4} TUMER K, AGOGINO A. Distributed agent-based air traffic flow management[C]//Proceedings of the 6th International Joint Conference on Autonomous Agents and Multiagent Systems (AAMAS 2007). Honolulu: ACM, 2007: 330-337.

% [17] 王迎 — 城市交通DRL硕士论文
\bibitem{ref13} 王迎. 城市交通信号控制深度强化学习算法研究[D]. 济南: 山东交通学院, 2024.

% [18] 赖建辉 — DRL交通控制硕士论文
\bibitem{ref16} 赖建辉. 基于深度强化学习的交通控制优化方法研究及实现[D]. 杭州: 浙江工业大学, 2020.

% [19] 张新武 — DRL交通信号优化硕士论文
\bibitem{ref17} 张新武. 基于深度强化学习算法的交通信号控制优化研究[D]. 太原: 太原科技大学, 2024.

% [20] 赵晓华等 — Q-learning+BP交叉口控制
\bibitem{ref14} 赵晓华, 石建军, 李振龙, 赵国勇. 基于Q-learning和BP神经元网络的交叉口信号灯控制[J]. 公路交通科技, 2007, 24(7): 99-102.

% [21] 承向军等 — Q-学习交通信号控制
\bibitem{ref15} 承向军, 常歆识, 杨肇夏. 基于Q-学习的交通信号控制方法[J]. 系统工程理论与实践, 2006(8): 136-140.

% [22] 卢守峰等 — 在线Q学习交叉口模型
\bibitem{ref18} 卢守峰, 张术, 刘喜敏. 平均排队长度差最小的单交叉口在线Q学习模型[J]. 公路交通科技, 2014, 31(11): 116-122.

% [23] 江波等 — DQN机场道口调速
\bibitem{ref19} 江波, 李彦冬, 刘威陇, 等. 应用深度Q学习的机场道口协同智能调速方法[J]. 航空计算技术, 2022, 52(1): 26-30.

% [24] 张春阳, 席雅琪 — AI交通控制综述
\bibitem{ref8} 张春阳, 席雅琪. 人工智能在城市交通控制中的应用综述[J]. 信息系统工程, 2024(07): 33-36.

% [25] Ghavamzadeh — 层级MARL (JAAMAS 2006)
\bibitem{ref12} GHAVAMZADEH M, MAHADEVAN S, MAKAR R. Hierarchical multi-agent reinforcement learning[J]. Autonomous Agents and Multi-Agent Systems, 2006, 13(2): 197-229.

% [26] Tumer & Agogino — 学习型多智能体ATM (IEEE IS 2009)
\bibitem{ref5} TUMER K, AGOGINO A. Improving air traffic management with a learning multiagent system[J]. IEEE Intelligent Systems, 2009, 24(1): 18-21.

% [27] Joby Aviation — S4技术指标
\bibitem{ref28} Joby Aviation. Joby S4 technical specifications[EB/OL]. 2024. \url{https://www.jobyaviation.com/}

% [28] Volocopter — VoloCity技术指标
\bibitem{ref29} Volocopter GmbH. VoloCity technical specifications[EB/OL]. 2024. \url{https://www.volocopter.com/}

% [29] Lilium — Lilium Jet技术指标
\bibitem{ref30} Lilium N.V. Lilium Jet technical specifications[EB/OL]. 2024. \url{https://www.lilium.com/}

% [30] GraphMixer — ICLR 2023 (用于区分命名)
\bibitem{ref25} CONG W, ZHANG S, CHEN J, et al. Do we really need complicated model architectures for temporal networks?[C]//Proceedings of the 11th International Conference on Learning Representations (ICLR 2023). Kigali: ICLR, 2023.

% [31] PTR-PPO — 优先轨迹回放PPO (arXiv 2021)
\bibitem{ref21} LIANG X, MA Y, FENG Y, et al. PTR-PPO: proximal policy optimization with prioritized trajectory replay[EB/OL]. arXiv preprint arXiv:2112.03798, 2021.

\end{thebibliography}

\end{document}
